% Options for packages loaded elsewhere
\PassOptionsToPackage{unicode}{hyperref}
\PassOptionsToPackage{hyphens}{url}
\PassOptionsToPackage{dvipsnames,svgnames,x11names}{xcolor}
\documentclass[
  10pt,
]{article}
\usepackage{xcolor}
\usepackage[a4paper,margin=2cm]{geometry}
\usepackage{amsmath,amssymb}
\setcounter{secnumdepth}{-\maxdimen} % remove section numbering
\usepackage{iftex}
\ifPDFTeX
  \usepackage[T1]{fontenc}
  \usepackage[utf8]{inputenc}
  \usepackage{textcomp} % provide euro and other symbols
\else % if luatex or xetex
  \usepackage{unicode-math} % this also loads fontspec
  \defaultfontfeatures{Scale=MatchLowercase}
  \defaultfontfeatures[\rmfamily]{Ligatures=TeX,Scale=1}
\fi
\usepackage{lmodern}
\ifPDFTeX\else
  % xetex/luatex font selection
\fi
% Use upquote if available, for straight quotes in verbatim environments
\IfFileExists{upquote.sty}{\usepackage{upquote}}{}
\IfFileExists{microtype.sty}{% use microtype if available
  \usepackage[]{microtype}
  \UseMicrotypeSet[protrusion]{basicmath} % disable protrusion for tt fonts
}{}
\makeatletter
\@ifundefined{KOMAClassName}{% if non-KOMA class
  \IfFileExists{parskip.sty}{%
    \usepackage{parskip}
  }{% else
    \setlength{\parindent}{0pt}
    \setlength{\parskip}{6pt plus 2pt minus 1pt}}
}{% if KOMA class
  \KOMAoptions{parskip=half}}
\makeatother
\usepackage{longtable,booktabs,array}
\newcounter{none} % for unnumbered tables
\usepackage{calc} % for calculating minipage widths
% Correct order of tables after \paragraph or \subparagraph
\usepackage{etoolbox}
\makeatletter
\patchcmd\longtable{\par}{\if@noskipsec\mbox{}\fi\par}{}{}
\makeatother
% Allow footnotes in longtable head/foot
\IfFileExists{footnotehyper.sty}{\usepackage{footnotehyper}}{\usepackage{footnote}}
\makesavenoteenv{longtable}
\setlength{\emergencystretch}{3em} % prevent overfull lines
\providecommand{\tightlist}{%
  \setlength{\itemsep}{0pt}\setlength{\parskip}{0pt}}
\usepackage{bookmark}
\IfFileExists{xurl.sty}{\usepackage{xurl}}{} % add URL line breaks if available
\urlstyle{same}
\hypersetup{
  pdftitle={Calibration of tabular foundation models versus logistic regression for clinical risk prediction: pre-registration of a simulation study and a real-data analysis},
  pdfauthor={Ebrahim Khaled Ebrahim},
  colorlinks=true,
  linkcolor={black},
  filecolor={Maroon},
  citecolor={Blue},
  urlcolor={black},
  pdfcreator={LaTeX via pandoc}}

\title{Calibration of tabular foundation models versus
logistic regression for clinical risk prediction:
pre-registration of a simulation study and a real-data
analysis}
\author{Ebrahim Khaled Ebrahim}
\date{2026-10-08}

\begin{document}
\maketitle

\textbf{Status.} Version 1.0, the final text for deposit.
The protocol is frozen at the moment it is registered on OSF
Registries. The registration must be made before the first
replicate of the main simulation run is written to disk; the
run scripts refuse to start until the registration is
recorded (Section 9.4). After the registration this file is
never edited. Every later change goes into
\texttt{DEVIATIONS.md}, dated, with its reason (Section 12).

\textbf{Author.} Ebrahim Khaled Ebrahim (sole author),
ebrahimkhaled@alexu.edu.eg, ORCID 0009-0006-7839-8778.

\textbf{Affiliation.} Department of Applied Statistics,
Faculty of Business, Alexandria University, Alexandria,
Egypt.

\textbf{Structure.} The simulation part follows the ADEMP
structure of Morris, White and Crowther (2019): Aims,
Data-generating mechanisms, Estimands, Methods, Performance
measures. The registration uses the OSF Registries
``Simulation Studies'' template, which follows the same
structure (Siepe et al.~2024).

\begin{center}\rule{0.5\linewidth}{0.5pt}\end{center}

\subsection{0. What was done and seen before this
registration}\label{what-was-done-and-seen-before-this-registration}

\subsubsection{0.1 Work before the
freeze}\label{work-before-the-freeze}

Before this text was frozen, the following work was done.
All of it used data that the main run never uses.

\begin{enumerate}
\def\labelenumi{\arabic{enumi}.}
\tightlist
\item
  \textbf{An earlier, unpublished pilot by the same author.}
  It used 3 replicates per scenario and small test sets. It
  is superseded, and no result from it is used.
\item
  \textbf{An oracle size study of the external tests}
  (\texttt{R/size\_check.R}).

  \begin{itemize}
  \tightlist
  \item
    2,000 replicates in each of 12 scenarios, with
    predictions equal to the true risk.
  \item
    It used its own coefficients and its own random stream,
    not the main design's.
  \item
    It is an auxiliary check (Section 5.3).
  \end{itemize}
\item
  \textbf{A Phase-1 pilot of this study} (2026-10-07):

  \begin{itemize}
  \tightlist
  \item
    every model arm on 6 scenarios x 20 replicates, then the
    full goodness-of-fit layer;
  \item
    TabPFN-3.5 on 5 scenarios x 20 replicates;
  \item
    timing and memory runs.
  \item
    The pilot \textbf{computed and inspected calibration
    summaries}: true slope and intercept, true ICI, and the
    rejection rates of the tests. They informed the changes
    listed in Section 0.2.
  \end{itemize}
\item
  \textbf{A read-only study of the external Stukel test}
  (Section 5.3), with its own random stream.
\item
  \textbf{A final smoke run} of every arm, the scoring and
  the test layer: 2 scenarios x 2 replicates (Section 9.2).
\item
  \textbf{A size study of the new Stukel row} of
  \texttt{ebrahim.gof} 2.10.0 at commit 19480aa (Section
  5.3): 4,000 nulls per cell, in the null design and with
  the seeds of item 4.
\item
  \textbf{Two engineering checks for the random-design tier}
  (Section 3.9), on the smoke stream:

  \begin{itemize}
  \tightlist
  \item
    a GPU memory and time probe of the three TFMs at 16 (n,
    p, predicted rows) cells on synthetic data, which sets
    the device rules;
  \item
    the classical arms and the test layer on 50 random
    designs (R0006-R0055), which give the per-fit costs used
    to choose D.
  \item
    Neither check computed or inspected a calibration
    summary.
  \end{itemize}
\item
  \textbf{A second smoke run} (Section 9.2): every arm on 2
  core scenarios x 2 replicates and on 5 random designs
  (R0001-R0005), then the scoring and the test layer at
  commit 19480aa.
\end{enumerate}

\textbf{Separate seed streams.}

\begin{itemize}
\tightlist
\item
  The pilot and the smoke runs drew every data set from
  \texttt{SeedSequence({[}scenario\_id,\ n\_index,\ rep,\ 777{]})}.
\item
  The main run draws from
  \texttt{SeedSequence({[}scenario\_id,\ n\_index,\ rep{]})}.
\item
  The random designs (Section 3.9) draw from a stream of
  their own,
  \texttt{SeedSequence({[}4242,\ 2,\ d,\ rep{]})}; their
  smoke data add 777.
\item
  The main run never touches the 777 stream.
\item
  Every main-run data set carries a SHA-256 fingerprint,
  which is asserted before every downstream step (Section
  3.8). A pilot file therefore cannot enter the main
  results.
\end{itemize}

\subsubsection{0.2 Changes made before the freeze, and
why}\label{changes-made-before-the-freeze-and-why}

Each row was decided by the author on 2026-10-07, after the
pilot. Rows marked ``pilot'' were informed by the pilot's
calibration summaries. They are design decisions made before
the main run, not deviations.

On 2026-10-08, before the freeze, the author changed the
draft's confirmatory hypotheses into an estimation-focused
protocol with pre-specified research questions and fixed
reporting rules (Section 7), by his own choice to describe
and estimate rather than test predictions, so that the
results determine the conclusions within rules fixed in
advance.

{\def\LTcaptype{none} % do not increment counter
\begin{longtable}[]{@{}
  >{\raggedright\arraybackslash}p{(\linewidth - 8\tabcolsep) * \real{0.1176}}
  >{\raggedright\arraybackslash}p{(\linewidth - 8\tabcolsep) * \real{0.1176}}
  >{\raggedright\arraybackslash}p{(\linewidth - 8\tabcolsep) * \real{0.2206}}
  >{\raggedright\arraybackslash}p{(\linewidth - 8\tabcolsep) * \real{0.2941}}
  >{\raggedright\arraybackslash}p{(\linewidth - 8\tabcolsep) * \real{0.2500}}@{}}
\toprule\noalign{}
\begin{minipage}[b]{\linewidth}\raggedright
\#
\end{minipage} & \begin{minipage}[b]{\linewidth}\raggedright
Item
\end{minipage} & \begin{minipage}[b]{\linewidth}\raggedright
Before
\end{minipage} & \begin{minipage}[b]{\linewidth}\raggedright
After (this registration)
\end{minipage} & \begin{minipage}[b]{\linewidth}\raggedright
Why
\end{minipage} \\
\midrule\noalign{}
\endhead
\bottomrule\noalign{}
\endlastfoot
C1 & Slope at small n (question Q1) & The plan of 2026-10-06
expected the median true slope b* of each TFM to exceed 1.1
at n \textless= 500 and pi = 10 \% in all four shapes (TFMs
under-confident) & No expected direction. Q1 estimates each
TFM's b* and its paired distance from 1 against FLIC on the
log scale, with the main cells at n \textless= 500 and pi =
10 \%, and labels the direction it finds by a fixed rule
(Section 7.1) & Pilot. In the two qualifying cells (S2 and
T3 at pi = 10 \%, n = 500), the TFM medians were 0.88-1.11
(MCSE 0.05-0.09), against FLIC's 0.84-0.86. At n = 200, pi =
3 \%, every classical arm was over-confident (FLIC 0.49) and
the TFMs were less so (0.71-1.00). These values informed the
emphasis of Q1: a paired comparison with FLIC, with both
directions labelled, in place of a single threshold of
1.1. \\
C2 & Rare events (question Q3) & The plan of 2026-10-06
expected each TFM to have a larger \textbar a*\textbar{}
than FLIC at pi = 3 \% at every n & No expected direction.
Q3 estimates the paired \textbar a*\textbar{} difference
against FLIC at pi = 3 \% at every n, shows the cells at n
\textgreater= 1000 first and each TFM's shipped softmax
temperature beside its result, and the temperature arms
estimate the temperature's effect directly (Section 7.3) &
Pilot. The sign depended on the temperature. At
S4\_pi03\_n3000, TabICLv2 and TabPFN v2 (both T = 0.9) had
a* = 0.29-0.32, against FLIC 0.02, Mitra-v2 0.00 and
TabPFN-3.5 0.02. At n = 200, no TFM had a larger
\textbar a*\textbar{} than FLIC. These values informed the
emphasis of Q3 on n \textgreater= 1000 and on the
temperature. \\
C3 & Primary TabPFN arm & ``TabPFN-3.5 if its licence is
accepted, else the newest ungated checkpoint'' & TabPFN-3.5
standard checkpoint, pinned by version and SHA-256 (Section
5.2) & The author accepted the licence, and the checkpoints
are in place. The choice between the standard and Fast
checkpoints is justified by timing in Section 5.2. \\
C4 & Temperature arms & TabPFN at T = 1.0 on 8 scenarios, R
= 200 & \texttt{tabpfn\_t1} = TabPFN-3.5 at T = 0.9. A new
arm \texttt{tabicl\_t1} = TabICLv2 at T = 1.0. Both on the
same 8 scenarios with R = 100. & Pilot. TabPFN-3.5 already
ships at T = 1.0, so the old arm would repeat the main arm.
TabICLv2 and TabPFN v2 ship at T = 0.9, which multiplies
every logit by 1/0.9 (measured slope 1.1105-1.1109) and
lowers the mean risk by 0.012-0.017. \\
C5 & LightGBM + Platt & Platt scaling as fitted & A Platt
slope \textless= 0 keeps the raw LightGBM scores (flag
\texttt{platt\_fallback}) & Pilot. At n = 200, pi = 3 \%, a
negative Platt slope reversed the ranking in 3 of 20
replicates. \\
C6 & Ridge-CV lambda path & glmnet's default path & 200
log-spaced values from glmnet's lambda\_max down to 1e-6 x
lambda\_max, with the early path stop switched off; an
optimum at either end is flagged & Pilot. The default path
stopped early, and the CV optimum sat at its lower end in 19
of 20 fits at S3\_pi30\_n1000. With the new path it was
interior in 49 of 50 smoke fits; the one exception is at the
upper end (the null model), and it is flagged. \\
C7 & Degenerate fits & Scored as ``ok'' & Labelled
``degenerate'', scored, kept and counted (Section 8.2) &
Pilot. At n = 200, pi = 3 \%, more than 100 predictions sat
at 0 or 1 in 12 of 20 GAM fits and 8 of 20 LR-MLE fits. \\
C8 & Mitra-v2 at T2, n = 3000 & Run like every scenario &
Recorded as ``does not fit in 6 GB GPU memory''. It is not
run there, and its context is never cut (Section 5.2). &
Pilot. 20 of 20 fits ran out of memory. A CPU fit had not
finished 100 rows after 15 minutes. \\
C9 & Device of TabPFN-3.5 & GPU, or CPU workers & Whole
scenario groups run on CPU workers, declared in Section 9.3.
T2\_pi10\_n3000 always runs on the CPU. & Pilot. TabPFN-3.5
runs out of GPU memory at T2, n = 3000. CPU and GPU
predictions agree to 2.6e-4. The schedule needs the CPU
workers to finish in about 3 days. \\
C10 & Replicates and test settings & TFMs R = 200
everywhere; legacy V2 R = 200; DeepGOF-1 B = 199 on real Y
and Y*; the battery on every classical replicate & Schedule
P4 (Section 6.3): TFMs R = 200 at n \textless= 1000 and R =
150 at n = 3000; classical R = 500; legacy V2 R = 100;
DeepGOF-1 B = 99 on real Y only, with the axes reading only
in T2; the battery on every arm for reps 1-R, and on the
oracle only for reps R+1 to 500 & Compute budget measured in
the pilot: the plan as written needed 4.0-4.7 days \\
C11 & Stukel (offset) & One of the battery tests & Its size
is reported beside its power, as for every test. The
exclusion from the detection ranking decided earlier the
same day is withdrawn by C21: the row is now the score test,
and it is ranked like the other tests (Section 5.3). & The
size study: 0.063 pooled at N = 3,334. The Stukel study
traced this to the likelihood-ratio statistic of its
upper-tail term, which the score form replaces. \\
C12 & The size study (review D5) & The 2,000-replicate
oracle study was the size study & The main size study is the
main run's own oracle nulls; the earlier study is auxiliary
(Section 5.3) & The earlier study used its own
coefficients \\
C13 & True ICI (review D6) & One loess estimator & The
pre-registered loess estimator stays primary.
\texttt{true\_ll\_ici} is reported beside it, and the oracle
row gives the loess floor (Section 4.1). & The loess of p*
on logit p-hat cannot follow the logistic curve: the oracle
gets 0.0008-0.0029, while \texttt{true\_ll\_ici} is 0 for
it \\
C14 & Class thresholds at pi = 3 \% (review D6) & Absolute
thresholds & The probability-scale thresholds are multiplied
by 0.3 at pi = 3 \% (Section 4.5) & At a mean risk of 0.03
the absolute thresholds were met by arms that failed
``weak'' \\
C15 & Few-event failures (review D8) & General pairing rule
& The pairing rule names the known few-event failures
(Section 8.2) & Ridge needs 3 events and LightGBM 2; the m =
100 ladder can draw 0 events \\
C16 & Mitra categoricals (review D9) & ``Declared
categorical to every TFM'' & Mitra-v2 receives integer codes
(Section 5.1) & Its interface has no categorical argument \\
C17 & e200 subsample (review D13) & Every event rate & pi =
30 \% only (Section 5.3) & It needs 200/pi rows: more than
N\_eval at 3 \% and 10 \% \\
C18 & Test package & ebrahim.gof 2.9.0 (+ localization from
2.10.0) & ebrahim.gof 2.10.0, which adds an ``Osius-Rojek
(external)'' row to \texttt{run.all.external()} & A
release-branch build that has every function used \\
C19 & LightGBM software & R \texttt{lightgbm} & Python
LightGBM 4.7.0, the same library, with the same grid and
early-stopping rule & R \texttt{lightgbm} is not in the
project library \\
C20 & GOF output files & csv & gzip-compressed csv
(\texttt{.csv.gz}); the reader handles both & Disk use \\
C21 & Package commit and the Stukel (offset) row &
\texttt{ebrahim.gof} 2.10.0 at commit 2c77f3c. Stukel
(offset) is the likelihood-ratio test, liberal at low event
rates (0.061-0.078). It is excluded from the detection
ranking: an exception in the size rule, and left out of the
detection comparisons, the detection curves and the ranking
in the summary tables. & 2.10.0 at commit \textbf{19480aa}.
Stukel (offset) is the \textbf{score test} of the same two
terms at the null (columns scaled to their information,
rank-aware inverse, df = 1 when no prediction reaches 0.5).
Its size is 0.035-0.053 in 4,000 nulls per cell (pi = 3, 10
and 30 \%; normal and t3 risks), with a p-value in every
replicate (Section 5.3). Stukel is in the size question
(Q4a) with no exception, and in the detection curves and the
summary tables, ranked like every other test. The test layer
refuses any build but 19480aa (Section 5.3). & The author
changed the package on 2026-10-07 after the Stukel study.
The likelihood ratio was liberal (0.061-0.078); the score
form holds its size. \\
C22 & Random-design robustness tier & None: 58 fixed
scenarios & A robustness tier of D = 1,000 randomly drawn
designs, one replicate each, with a pre-specified meta-model
and a transfer check (Sections 3.9, 6.4 and 7.6) & The
conclusions of a fixed factorial can be challenged as
depending on the chosen settings. The tier draws every
setting at random and checks whether the conclusions carry
over. \\
C23 & Outputs for reuse & Per-replicate predictions
deposited & (a) the per-replicate predictions and test
results released as an open benchmark, with a loader and a
data dictionary; (b) a calibration audit in one call, in R
and from Python; (c) a minimum-n-for-calibration table by
event rate for every model, as a primary output; (d) a
TRIPOD+AI alignment table. TabPFN weights are never archived
(Section 14). & Reuse by practitioners and methodologists \\
\end{longtable}
}

\begin{center}\rule{0.5\linewidth}{0.5pt}\end{center}

\subsection{1. Background and
rationale}\label{background-and-rationale}

Tabular foundation models (TFMs) such as TabPFN, Mitra and
TabICL predict a binary outcome in context. They read the
development sample as a prompt and return a probability for
each new row, with no gradient fitting on the user's data.
They are being tried for clinical risk prediction. There,
calibration (agreement between predicted and actual risk)
decides whether a prediction can be used for a clinical
decision.

Recent studies report TFM calibration with the Brier score,
log-loss, expected calibration error, conformal coverage or
a calibration plot read by eye. None of them does the
following five things, and this study does:

\begin{enumerate}
\def\labelenumi{\arabic{enumi}.}
\tightlist
\item
  score TFM predictions against the true risk;
\item
  check the size of the external-validation goodness-of-fit
  tests on the TFMs' own prediction distributions;
\item
  localize the miscalibration in the risk range or in
  covariate subgroups;
\item
  place TFMs on the calibration hierarchy of Van Calster et
  al.~(2016) and on the sample-size criteria of Riley et
  al.~(2020);
\item
  evaluate a test-guided recalibration rule.
\end{enumerate}

Two further parts make the results easier to trust and to
reuse. A tier of randomly drawn designs checks whether the
conclusions depend on the chosen scenarios (Section 3.9).
The predictions, a minimum-n table and a one-call
calibration audit are released for others to use (Section
14).

\textbf{Prior information, disclosed.} Section 0 lists
everything that was run and seen before this registration,
and which design choices it informed. The pilot also
informed the emphasis of research questions Q1 and Q3
(Section 7). The protocol predicts the direction of no
result.

\begin{center}\rule{0.5\linewidth}{0.5pt}\end{center}

\subsection{2. Aims}\label{aims}

\textbf{A1.} Quantify how far TFM risk predictions are from
the true risk, as a function of development sample size and
event rate. The comparators are Firth logistic regression
with intercept correction (FLIC), maximum-likelihood
logistic regression, ridge regression, a generalized
additive model and LightGBM.

\textbf{A2.} Find, for each model arm, the smallest
development size at which it reaches pre-specified
calibration targets (``minimum n for calibration''), and
compare it with Riley's minimum sample size.

\textbf{A3.} Measure the size (type I error) and the
detection rate of the external-validation goodness-of-fit
tests when they are applied to each arm's frozen
predictions.

\textbf{A4.} Measure how often the localization procedure
names the correct source of miscalibration when a known
covariate interaction is missed.

\textbf{A5.} Measure how much a small local sample repairs
calibration (a recalibration ladder), and whether a
test-guided rule picks the right level.

\textbf{A6.} Check on real clinical data (Section 10)
whether the patterns found in A1-A5 appear in real patients.

\textbf{A7.} Check whether the conclusions of A1 and A3 hold
across the design space, on randomly drawn designs (Section
3.9), and map the regions where each TFM is better or worse
calibrated than FLIC.

\textbf{A8.} Give practitioners and methodologists outputs
they can reuse (Section 14):

\begin{itemize}
\tightlist
\item
  a minimum-n-for-calibration table by event rate for each
  model (a primary output, Section 11);
\item
  a calibration audit in one call for any model's
  predictions;
\item
  the full per-replicate predictions and test results as an
  open benchmark;
\item
  an alignment of the reporting with TRIPOD+AI.
\end{itemize}

\begin{center}\rule{0.5\linewidth}{0.5pt}\end{center}

\subsection{3. Data-generating mechanisms
(simulation)}\label{data-generating-mechanisms-simulation}

\subsubsection{3.1 Covariates}\label{covariates}

The core design has eight covariates:

\begin{itemize}
\tightlist
\item
  X1, \ldots, X6 continuous, marginally N(0, 1); corr(X1,
  X2) = 0.5; all other pairs independent;
\item
  X7 \textasciitilde{} Bernoulli(0.5);
\item
  X8 \textasciitilde{} Bernoulli(0.3).
\end{itemize}

\subsubsection{3.2 Outcome models
(shapes)}\label{outcome-models-shapes}

Let
\texttt{L\ =\ 1.0\ X1\ +\ 0.8\ X2\ +\ 0.6\ X3\ +\ 0.4\ X4\ +\ 0.2\ X5\ +\ 0.0\ X6\ +\ 0.5\ X7\ +\ 0.3\ X8}
be the base linear term (X6 is a pure noise covariate).
Then:

{\def\LTcaptype{none} % do not increment counter
\begin{longtable}[]{@{}
  >{\raggedright\arraybackslash}p{(\linewidth - 6\tabcolsep) * \real{0.1667}}
  >{\raggedright\arraybackslash}p{(\linewidth - 6\tabcolsep) * \real{0.3167}}
  >{\raggedright\arraybackslash}p{(\linewidth - 6\tabcolsep) * \real{0.2500}}
  >{\raggedright\arraybackslash}p{(\linewidth - 6\tabcolsep) * \real{0.2667}}@{}}
\toprule\noalign{}
\begin{minipage}[b]{\linewidth}\raggedright
Shape
\end{minipage} & \begin{minipage}[b]{\linewidth}\raggedright
True linear predictor eta
\end{minipage} & \begin{minipage}[b]{\linewidth}\raggedright
Link
\end{minipage} & \begin{minipage}[b]{\linewidth}\raggedright
Correctly specified arms
\end{minipage} \\
\midrule\noalign{}
\endhead
\bottomrule\noalign{}
\endlastfoot
S1 linear & alpha + kappa L & logit & LR-MLE, FLIC, Ridge,
GAM \\
S2 smooth non-linear & alpha + kappa {[}L + gamma ((X1\^{}2
- 1)/sqrt(2) + (X3\^{}3 - 3 X3)/sqrt(6)){]}, gamma = 0.5 &
logit & GAM \\
S3 interactions & alpha + kappa {[}L + 0.5 X1 X7 + 0.5 X3
X4{]} & logit & none of the statistical arms as fitted (main
effects only) \\
S4 asymmetric link & alpha + kappa L & complementary
log-log: p* = 1 - exp(-exp(eta)) & none \\
\end{longtable}
}

The true risk of a row is p* = logit\^{}\{-1\}(eta) in
S1-S3, and 1 - exp(-exp(eta)) in S4. The outcome is Y
\textasciitilde{} Bernoulli(p*).

In S3, the interaction X1 x X7 defines the \emph{true
miscalibrated subgroup}. Patients with X7 = 1 have a steeper
risk gradient in X1 than any main-effects model can
represent. This subgroup is the target of the localization
analysis (Sections 5.4 and 11).

These formulas are implemented in \texttt{code/dgp.py}. An
independent re-implementation from this text agreed with it
to 5.6e-16 on all 58 scenarios.

\subsubsection{3.3 Signal strength and event
rate}\label{signal-strength-and-event-rate}

\begin{itemize}
\tightlist
\item
  For each shape (and each severity gamma in the stress
  set), the scale factor kappa is solved numerically on
  10\^{}6 fixed draws, so that the oracle AUC (the AUC of
  p*) is 0.80 at an event rate of 10 \%.
\item
  For each shape and each event rate pi in \{3 \%, 10 \%, 30
  \%\}, the intercept alpha is then solved on the same draws
  so that E{[}p*{]} = pi.
\item
  The realized prevalence equals pi to 2e-15 in every
  scenario.
\item
  Because kappa is fixed at 10 \%, the oracle AUC is
  0.77-0.83 at 3 \% and 30 \%.
\item
  The solved values and the realized oracle AUC of every
  scenario are listed in Appendix A.
\end{itemize}

\subsubsection{3.4 Development sample
size}\label{development-sample-size}

n in \{200, 500, 1000, 3000\}. Each scenario is also
reported in two ways:

\begin{itemize}
\tightlist
\item
  as events per variable (EPV = n pi / p);
\item
  against Riley et al.'s (2020) minimum sample size for a
  binary outcome, computed with the R package
  \texttt{pmsampsize} (expected shrinkage 0.9, Cox-Snell
  R\^{}2 from the oracle AUC, p parameters).
\end{itemize}

The scenario n = 200 at pi = 3 \% (about 6 events, EPV 0.75)
is kept on purpose: it is the regime where TFMs might help.
Degenerate fits there are an outcome, not a reason to drop
the scenario.

\subsubsection{3.5 Core factorial}\label{core-factorial}

4 shapes x 3 event rates x 4 sizes = \textbf{48 core
scenarios} (ids 1-48).

\subsubsection{3.6 Stress set}\label{stress-set}

At pi = 10 \% and n in \{500, 3000\}, five stress types (10
scenarios, ids 49-58):

{\def\LTcaptype{none} % do not increment counter
\begin{longtable}[]{@{}
  >{\raggedright\arraybackslash}p{(\linewidth - 4\tabcolsep) * \real{0.1250}}
  >{\raggedright\arraybackslash}p{(\linewidth - 4\tabcolsep) * \real{0.7500}}
  >{\raggedright\arraybackslash}p{(\linewidth - 4\tabcolsep) * \real{0.1250}}@{}}
\toprule\noalign{}
\begin{minipage}[b]{\linewidth}\raggedright
ID
\end{minipage} & \begin{minipage}[b]{\linewidth}\raggedright
Change from the core design
\end{minipage} & \begin{minipage}[b]{\linewidth}\raggedright
Shape
\end{minipage} \\
\midrule\noalign{}
\endhead
\bottomrule\noalign{}
\endlastfoot
T1 & X1-X6 drawn as Student t with 3 df through a Gaussian
copula (X1 and X2 keep correlation 0.5), scaled to unit
variance & S1 \\
T2 & p = 30: the 8 core covariates plus 22 independent N(0,
1) noise columns & S1 \\
T3 & X8 replaced by a 6-level categorical with equal level
probabilities and level effects (0, 0.3, -0.3, 0.6, -0.6,
0.9) inside L. It is declared categorical to TabPFN and
TabICL, given as integer codes to Mitra-v2, and dummy-coded
for the statistical arms. & S3 \\
T4a & S2 with severity gamma = 0.25 & S2 \\
T4b & S2 with severity gamma = 1.0 & S2 \\
\end{longtable}
}

Together with S2 at gamma = 0.5, T4a and T4b form a severity
ladder for detection-power curves.

\textbf{Total: 58 scenarios.} Each has a fixed integer
\texttt{scenario\_id} (Appendix A).

\subsubsection{3.7 Data drawn per
replicate}\label{data-drawn-per-replicate}

Each replicate draws the following sets, independently and
fresh:

{\def\LTcaptype{none} % do not increment counter
\begin{longtable}[]{@{}
  >{\raggedright\arraybackslash}p{(\linewidth - 4\tabcolsep) * \real{0.1667}}
  >{\raggedright\arraybackslash}p{(\linewidth - 4\tabcolsep) * \real{0.3667}}
  >{\raggedright\arraybackslash}p{(\linewidth - 4\tabcolsep) * \real{0.4667}}@{}}
\toprule\noalign{}
\begin{minipage}[b]{\linewidth}\raggedright
Set
\end{minipage} & \begin{minipage}[b]{\linewidth}\raggedright
Size
\end{minipage} & \begin{minipage}[b]{\linewidth}\raggedright
Use
\end{minipage} \\
\midrule\noalign{}
\endhead
\bottomrule\noalign{}
\endlastfoot
Development & n & fitting, or the context of a TFM \\
Evaluation & N\_eval = max(1000, ceil(100/pi)): 3,334 at 3
\%, 1,000 at 10 \% and 30 \% & oracle estimands (p* known)
and the test battery on fresh Y; at least 100 expected
events (Riley et al.~2021) \\
Recalibration & 500 & recalibration ladder; the m = 100
level uses the first 100 rows \\
\end{longtable}
}

A \textbf{probe set} of 500 rows per scenario is drawn once,
from \texttt{SeedSequence({[}scenario\_id,\ 99999{]})}, and
is the same in every replicate. It is used for the
prediction-instability plots.

Smaller validation samples are nested subsamples of the
stored evaluation rows (the first ceil(50/pi) and
ceil(200/pi) rows); no model is refitted. The e50 subsample
exists at every event rate. The e200 subsample exists only
at pi = 30 \%, because at 3 \% and 10 \% it would need more
rows than N\_eval.

\subsubsection{3.8 Random numbers and
fingerprints}\label{random-numbers-and-fingerprints}

\begin{itemize}
\tightlist
\item
  \textbf{Data.}

  \begin{itemize}
  \tightlist
  \item
    Every data set is drawn in Python from
    \texttt{numpy.random.SeedSequence({[}scenario\_id,\ n\_index,\ rep{]})},
    with rep \textgreater= 1.
  \item
    The seed is spawned into independent child streams for
    the development, evaluation and recalibration draws.
  \item
    Model seeds and null-outcome seeds come from the same
    entropy with fixed arm codes.
  \item
    The pilot and smoke stream adds a fourth entry, 777
    (Section 0.1).
  \end{itemize}
\item
  \textbf{Fingerprints.}

  \begin{itemize}
  \tightlist
  \item
    A SHA-256 fingerprint of the bytes of (X, p*, Y) of
    every set is stored with the replicate.
  \item
    Every downstream step recomputes and asserts it before
    using the data: adding an arm, scoring, and the R
    layer's cross-file checks. A mismatch stops the run.
  \item
    Two replicates with identical data in one run stop the
    run, which would catch a hidden reset of the random
    stream.
  \end{itemize}
\item
  \textbf{R procedures.}

  \begin{itemize}
  \tightlist
  \item
    Procedures that draw random numbers (the Monte Carlo
    references of \texttt{localize.external()} and
    \texttt{deepgof1.external()}, glmnet fold assignment,
    null outcomes drawn in R) receive an explicit integer
    seed. The seed is derived from a string naming the
    scenario, replicate, arm and test.
  \item
    The generator is reset to Mersenne-Twister before each
    such call.
  \item
    Equal outcome fingerprints across replicates of a
    scenario are reported as an error.
  \end{itemize}
\end{itemize}

\subsubsection{3.9 Random-design robustness
tier}\label{random-design-robustness-tier}

\textbf{Purpose.} The 58 scenarios fix every factor at a few
levels. A reader may ask whether the conclusions hold away
from those levels. In this tier every design draws its own
setting, so that the conclusions can be checked across the
whole design space, and the regions where they fail can be
mapped (Sections 6.4 and 7.6).

\textbf{Design.} D = 1,000 designs
(\texttt{R0001}-\texttt{R1000}), each with \textbf{one
replicate}. Design d draws its setting from
\texttt{SeedSequence({[}4242,\ 0,\ d{]})}:

{\def\LTcaptype{none} % do not increment counter
\begin{longtable}[]{@{}
  >{\raggedright\arraybackslash}p{(\linewidth - 2\tabcolsep) * \real{0.6000}}
  >{\raggedright\arraybackslash}p{(\linewidth - 2\tabcolsep) * \real{0.4000}}@{}}
\toprule\noalign{}
\begin{minipage}[b]{\linewidth}\raggedright
Factor
\end{minipage} & \begin{minipage}[b]{\linewidth}\raggedright
Distribution
\end{minipage} \\
\midrule\noalign{}
\endhead
\bottomrule\noalign{}
\endlastfoot
event rate pi & log-uniform on {[}0.02, 0.40{]} \\
development size n & log-uniform on {[}150, 4000{]},
rounded \\
number of predictors p & integer uniform on {[}4, 30{]} \\
predictor correlation rho & uniform on {[}0, 0.6{]} \\
oracle AUC target & uniform on {[}0.65, 0.90{]} \\
severity of non-linearity w\_nl, of interaction w\_int, of
link asymmetry w\_link & each uniform on {[}0, 1{]} \\
fraction of binary or categorical predictors & uniform on
{[}0, 0.5{]} \\
\end{longtable}
}

\textbf{Covariates.}

\begin{itemize}
\tightlist
\item
  A latent Gaussian vector with exchangeable correlation
  rho.
\item
  Of the p predictors, round(fraction x p) are discrete and
  the rest are continuous (the latent values).
\item
  Each discrete predictor is categorical with probability
  0.25 (3-6 levels, equal probabilities, cut at normal
  quantiles of its latent value); otherwise it is binary,
  1\{z \textgreater{} q\}, with a prevalence drawn on
  {[}0.2, 0.5{]}.
\item
  Categorical predictors are declared to the TFMs and coded
  for the statistical arms as in T3 (Section 5.1).
\end{itemize}

\textbf{Outcome.} The shape is a random mixture built from
the core design's own shape generators
(\texttt{code/dgp.py}):

\begin{itemize}
\tightlist
\item
  The linear term L has coefficients drawn per design: a
  random sign, a size uniform on {[}0.2, 1{]}, and each
  predictor after the second is pure noise with probability
  0.25; categorical levels get N(0, 0.5\^{}2) effects. L is
  rescaled to the SD of the core design's L (1.755), so that
  the weights below mean what gamma means in the core
  design.
\item
  N is S2's non-linear term at gamma = 1, and I is S3's
  interaction term (Section 3.2), both evaluated by
  \texttt{dgp.shape\_score} on the design's predictors in
  the core roles: x1 = the first continuous predictor, x3 =
  the second, x4 = the third (the first if there are only
  two), x7 = the first binary predictor (1\{z \textgreater{}
  0\} of the last continuous predictor if there is none).
\item
  eta = alpha + kappa {[}L + w\_nl N + w\_int I{]}, and p* =
  (1 - w\_link) logit\^{}\{-1\}(eta) + w\_link (1 -
  exp(-exp(eta))): w\_link = 1 is S4's complementary log-log
  link.
\item
  kappa is solved so that the oracle AUC equals the design's
  target at its own event rate, and alpha so that E{[}p*{]}
  = pi, on 200,000 calibration draws from
  \texttt{SeedSequence({[}4242,\ 1,\ d{]})}. Every design
  was solved: the realized prevalence equals pi to 6e-14,
  and the oracle AUC equals its target to 2e-10.
\end{itemize}

\textbf{Data per design.}

\begin{itemize}
\tightlist
\item
  Development set of size n.
\item
  Evaluation set sized as in the core design: N\_eval =
  max(1000, ceil(100/pi)), so at least 100 expected events
  (1,000-5,000 rows).
\item
  No recalibration or probe set: there is one replicate and
  no recalibration ladder.
\item
  The data come from
  \texttt{SeedSequence({[}4242,\ 2,\ d,\ 1{]})}, a stream
  that the core design and the pilot never use, and they
  carry the same SHA-256 fingerprints as the core design
  (Section 3.8).
\end{itemize}

\textbf{Arms.} Oracle, LR-MLE, FLIC, Ridge-CV, GAM, LightGBM
+ Platt (with its raw column), \textbf{TabPFN-3.5 (the main
arm)}, TabICLv2 and Mitra-v2, with the settings of Section
5.2. There are no legacy or temperature arms and no
recalibration ladder.

\textbf{Device rules} (measured on the 6 GB card with the
run's 0.8 memory cap: pilot Section 2 and
\texttt{results/random\_memory\_probe.csv}):

\begin{itemize}
\tightlist
\item
  \textbf{TabPFN-3.5} holds the development rows and the
  predicted rows together, so its size is (n + N\_eval) x p,
  not n x p alone: a 1,500 x 20 design with 5,000 predicted
  rows ran out of memory, while 4,000 x 8 with 2,000
  predicted rows fitted. Every measured cell up to 84,000
  fitted (peak 4.4 GB), and every cell from 99,000 up ran
  out of memory. A design runs TabPFN-3.5 on the GPU when (n
  + N\_eval) x p \textless= 80,000, and on a CPU worker
  otherwise (179 of the 1,000 designs). The device is
  recorded per design; CPU and GPU predictions agree to
  2.6e-4 (Section 5.2).
\item
  \textbf{Mitra-v2} predicts in batches, and its memory
  grows with n\^{}2 x p: every measured cell up to 8.0e7
  fitted (2,000 x 20 at 4.1 GB, 3,000 x 8 at 3.7 GB), and
  every cell from 9.4e7 up failed (2,500 x 15, 3,500 x 8,
  4,000 x 6, and T2 at 3,000 x 30). Above n\^{}2 x p = 8.0e7
  Mitra-v2 is \textbf{recorded as ``does not fit in 6 GB GPU
  memory''}, by the rule used for T2\_pi10\_n3000 (170 of
  the 1,000 designs). It is not run there, and its context
  is never cut. Its meta-model covers only the designs where
  it ran, and that region is reported.
\end{itemize}

\textbf{Test layer.} As in the core design (Section 5.3), at
B = 99, with one change of reading: DeepGOF-1 uses the axes
reading when the design has more than 12 predictors, and the
combined reading otherwise. In the core design this rule
picks out exactly T2.

\textbf{Files.} The solved designs are in
\texttt{code/random\_designs.csv}
(\texttt{python\ code/dgp\_random.py\ build\ -\/-D\ 1000};
\texttt{check} re-solves them and compares). Appendix B
summarizes them.

\begin{center}\rule{0.5\linewidth}{0.5pt}\end{center}

\subsection{4. Estimands}\label{estimands}

All oracle estimands are computed on the evaluation set of
each replicate, where the true risk p* of every row is
known. Predictions p-hat are clipped to {[}1e-6, 1 - 1e-6{]}
for every computation. The number of predictions within 1e-6
of 0 or 1 is recorded per arm and replicate.

\subsubsection{4.1 Primary estimands
(oracle)}\label{primary-estimands-oracle}

\begin{enumerate}
\def\labelenumi{\arabic{enumi}.}
\tightlist
\item
  \textbf{True calibration slope b*.} The slope of a
  quasi-binomial generalized linear model of p* on
  logit(p-hat), with a logit link.

  \begin{itemize}
  \tightlist
  \item
    b* \textgreater{} 1 means under-confidence (predictions
    too close to the mean).
  \item
    b* \textless{} 1 means over-confidence.
  \end{itemize}
\item
  \textbf{True calibration intercept a*.} The intercept of a
  quasi-binomial model of p* with offset logit(p-hat) and no
  slope (calibration-in-the-large on the logit scale).
\item
  \textbf{True integrated calibration index (true ICI).} The
  mean over evaluation rows of \textbar p-hat -
  c(p-hat)\textbar.

  \begin{itemize}
  \tightlist
  \item
    c(.) is a loess fit (degree 2, span 0.75, Gaussian
    family, no robustness iterations, exact fit as R's
    \texttt{surface\ =\ "direct"}) of \textbf{p*} (not Y) on
    logit(p-hat), evaluated at the evaluation rows.
  \item
    This loess cannot follow the logistic curve exactly, so
    the estimator has a floor: 0.0008-0.0029 for the oracle
    in the pilot.
  \item
    The oracle row is reported in every table as that floor.
  \item
    The \textbf{log-scale version \texttt{true\_ll\_ici}} is
    reported beside it. It smooths logit p* on logit p-hat
    and maps back. It is exactly 0 for the oracle, and it
    equals the OCE whenever logit p* is linear in logit
    p-hat.
  \item
    The true ICI remains the primary estimand.
  \end{itemize}
\end{enumerate}

\subsubsection{4.2 Further oracle estimands
(secondary)}\label{further-oracle-estimands-secondary}

\begin{itemize}
\tightlist
\item
  True E50, E90 and E\_max: the median, 90th percentile and
  maximum of \textbar p-hat - c(p-hat)\textbar.
\item
  \textbf{Oracle calibration error (OCE)} = mean
  \textbar p-hat - p*\textbar. It is the individual-level
  error. It separates ``miscalibrated'' from ``calibrated on
  average but poorly resolved''.
\item
  Mean p-hat minus mean p* (calibration-in-the-large on the
  probability scale).
\item
  True discrimination: the expected AUC of p-hat against p*,
  and the oracle AUC.
\end{itemize}

\subsubsection{4.3 Empirical estimands
(secondary)}\label{empirical-estimands-secondary}

These are computed with the observed Y on the evaluation
set, as a real validation study would:

\begin{itemize}
\tightlist
\item
  empirical ICI, E50, E90 and E\_max (Austin and Steyerberg
  2019: loess of Y on logit(p-hat), degree 2, no robustness
  iterations);
\item
  empirical calibration slope (joint fit of Y on
  logit(p-hat)) and calibration-in-the-large (offset fit);
\item
  Brier score, log-loss and AUC;
\item
  10-bin expected calibration error (reported only for
  comparison with the machine-learning literature);
\item
  the gap between empirical and true ICI per replicate (how
  wrong a typical validation estimate is).
\end{itemize}

\subsubsection{4.4 Prediction instability
(secondary)}\label{prediction-instability-secondary}

On the fixed probe set: the per-patient spread of p-hat
across replicates (standard deviation and 2.5-97.5 \% range)
(Riley and Collins 2023).

\subsubsection{4.5 Calibration class of each
replicate}\label{calibration-class-of-each-replicate}

{\def\LTcaptype{none} % do not increment counter
\begin{longtable}[]{@{}
  >{\raggedright\arraybackslash}p{(\linewidth - 4\tabcolsep) * \real{0.2667}}
  >{\raggedright\arraybackslash}p{(\linewidth - 4\tabcolsep) * \real{0.4500}}
  >{\raggedright\arraybackslash}p{(\linewidth - 4\tabcolsep) * \real{0.2833}}@{}}
\toprule\noalign{}
\begin{minipage}[b]{\linewidth}\raggedright
Class
\end{minipage} & \begin{minipage}[b]{\linewidth}\raggedright
Rule (pi = 10 \% and 30 \%)
\end{minipage} & \begin{minipage}[b]{\linewidth}\raggedright
Rule at pi = 3 \%
\end{minipage} \\
\midrule\noalign{}
\endhead
\bottomrule\noalign{}
\endlastfoot
mean & abs(mean p-hat - mean p*) \textless= 0.01 &
\textless= 0.003 \\
weak & abs(a*) \textless= 0.1 and abs(b* - 1) \textless= 0.1
& the same \\
moderate & true ICI \textless= 0.015 & \textless= 0.0045 \\
strong-approx & true E90 \textless= 0.03 & \textless=
0.009 \\
\end{longtable}
}

\begin{itemize}
\tightlist
\item
  \textbf{Relative thresholds at 3 \%.} At pi = 3 \% the
  three probability-scale thresholds are multiplied by pi /
  0.10 = 0.3, because the absolute values are easy to meet
  when the mean risk is 0.03. The weak rule is on the logit
  scale and is not changed.
\item
  \textbf{Nesting.} The classes are nested, as in the
  hierarchy of Van Calster et al.~(2016). An (arm,
  replicate) is placed in the highest class whose rule holds
  \emph{together with every lower rule}. Each rule is also
  stored on its own.
\item
  \textbf{Sensitivity.} The class-based results are repeated
  at half and at double every threshold, and with
  \texttt{true\_ll\_ici} in place of the true ICI.
\end{itemize}

\textbf{Minimum n for calibration} of an arm in a (shape,
pi, class) is the smallest n in \{200, 500, 1000, 3000\} at
which the arm meets the class in at least 80 \% of
replicates. If no n qualifies, the value is ``\textgreater{}
3000''.

\begin{center}\rule{0.5\linewidth}{0.5pt}\end{center}

\subsection{5. Methods (model
arms)}\label{methods-model-arms}

\subsubsection{5.1 Common rules}\label{common-rules}

\begin{itemize}
\tightlist
\item
  \textbf{Paired design.} Every arm is fitted on the
  \emph{same} development set of a replicate. It predicts
  the same evaluation, recalibration and probe rows.
\item
  \textbf{No tuning on evaluation data.} No arm is tuned on
  evaluation or recalibration data. The TFMs are used with
  their shipped settings and are not tuned at all. The only
  exception is the softmax temperature of the two
  temperature arms.
\item
  \textbf{Categorical covariates.} The binary covariates and
  the 6-level factor of T3 are declared categorical to
  TabPFN (\texttt{categorical\_features\_indices}) and to
  TabICLv2 (pandas categories). \textbf{Mitra-v2 receives
  them as integer codes}, because its interface has no
  categorical argument; this is how AutoGluon itself passes
  categoricals to Mitra.
\item
  \textbf{Provenance.} Every version, checkpoint file,
  revision hash and SHA-256 is recorded per replicate and in
  \texttt{results/environment.json} (Section 13).
\end{itemize}

\subsubsection{5.2 Arms}\label{arms}

{\def\LTcaptype{none} % do not increment counter
\begin{longtable}[]{@{}
  >{\raggedright\arraybackslash}p{(\linewidth - 6\tabcolsep) * \real{0.1500}}
  >{\raggedright\arraybackslash}p{(\linewidth - 6\tabcolsep) * \real{0.2333}}
  >{\raggedright\arraybackslash}p{(\linewidth - 6\tabcolsep) * \real{0.1333}}
  >{\raggedright\arraybackslash}p{(\linewidth - 6\tabcolsep) * \real{0.4833}}@{}}
\toprule\noalign{}
\begin{minipage}[b]{\linewidth}\raggedright
Arm (code name)
\end{minipage} & \begin{minipage}[b]{\linewidth}\raggedright
Implementation
\end{minipage} & \begin{minipage}[b]{\linewidth}\raggedright
Role
\end{minipage} & \begin{minipage}[b]{\linewidth}\raggedright
Settings
\end{minipage} \\
\midrule\noalign{}
\endhead
\bottomrule\noalign{}
\endlastfoot
Oracle (\texttt{oracle}) & p* itself & defines the null and
the truth & the ``perfect model'' row; the main size study
(Section 5.3) \\
LR-MLE (\texttt{lr}) & R \texttt{stats::glm}, binomial,
logit, main effects & classical reference & separation is
flagged, not hidden \\
\textbf{FLIC} (\texttt{flic}) & R
\texttt{logistf::logistf(...,\ flic\ =\ TRUE,\ pl\ =\ FALSE)}
(Puhr et al.~2017), main effects & \textbf{primary
statistical reference} & package defaults \\
Ridge-CV (\texttt{ridge}) & R \texttt{glmnet}, alpha = 0,
lambda = \texttt{lambda.min} by 10-fold CV deviance, folds
stratified on Y & penalization arm & \textbf{lambda path:}
200 log-spaced values from glmnet's own lambda\_max down to
1e-6 x lambda\_max, with glmnet's early path stop
(\texttt{fdev}) switched off for this fit;
\texttt{lambda\_at\_edge} flags an optimum at either end \\
GAM (\texttt{gam}) & R \texttt{mgcv::gam}, REML,
\texttt{s(x,\ k\ =\ 5)} for each continuous covariate,
linear binary and factor terms & flexible statistical arm &
correctly specified for S2 \\
LightGBM + Platt (\texttt{lgbm}), and raw
(\texttt{lgbm\_raw}) & Python LightGBM 4.7.0 &
machine-learning reference & 5-fold stratified CV over
num\_leaves in \{7, 15, 31\} x min\_data\_in\_leaf in \{10,
30\}, learning rate 0.05, up to 2,000 rounds with early
stopping after 50; refit on all n at the CV-best setting and
round count. \textbf{Platt scaling} is a logistic fit of Y
on the logit of the out-of-fold scores, applied to the
refit. \textbf{If the Platt slope is \textless= 0, the arm
keeps the raw LightGBM scores and records
\texttt{platt\_fallback}}. A Platt fit that does not
converge is a failure. Isotonic scaling appears only in the
ladder. \\
\textbf{TabPFN-3.5} (\texttt{tabpfn}) & \texttt{tabpfn}
9.1.0, \texttt{ModelVersion.V3\_5}, checkpoint
tabpfn-v3.5-20260909 (file, repository, revision and SHA-256
in Section 13) & primary TFM & shipped settings:
\texttt{n\_estimators} ``auto'' and softmax temperature
\textbf{1.0}, both recorded per replicate; pinned by version
(\texttt{V3\_5}) and SHA-256 \\
TabPFN-3.5 at T = 0.9 (\texttt{tabpfn\_t1}) & the same
checkpoint, \texttt{softmax\_temperature\ =\ 0.9}, the same
seed as \texttt{tabpfn} & temperature arm & 8 scenarios
(S1-S4 at pi = 10 \%, n in \{200, 3000\}), R = 100 \\
Real-TabPFN-2, legacy (\texttt{tabpfn\_legacy}) &
\texttt{tabpfn} 9.1.0, \texttt{ModelVersion.V2}, its default
classifier checkpoint (file, repository and SHA-256 in
Section 13) & link to the earlier TabPFN literature &
shipped settings (8 estimators, T = 0.9); the 16 core
scenarios at pi = 10 \%, R = 100 \\
\textbf{TabICLv2} (\texttt{tabicl}) & \texttt{tabicl} 2.2.0,
checkpoint tabicl-classifier-v2-20260212 (file, repository,
revision and SHA-256 in Section 13) & third TFM family &
shipped settings (8 estimators, softmax temperature
\textbf{0.9}). Its documented lower bound is 300 training
rows, so results at n = 200 are reported and flagged. \\
TabICLv2 at T = 1.0 (\texttt{tabicl\_t1}) & the same
checkpoint, \texttt{softmax\_temperature\ =\ 1.0}, the same
seed as \texttt{tabicl} & temperature arm & the same 8
scenarios, R = 100 \\
\textbf{Mitra-v2} (\texttt{mitra}) &
\texttt{MitraClassifier} from \texttt{autogluon.tabular}
1.6.3, weights autogluon/mitra-classifier-2 (revision and
SHA-256 in Section 13) & second TFM family & in-context only
(no fine-tuning, \texttt{n\_estimators\ =\ 1}), with the
whole development set as context. A 16-row validation set
taken from the development rows is passed only so that
AutoGluon holds nothing out; it does not change the
predictions. See below for precision, sign flips and the T2
cell. \\
\end{longtable}
}

\textbf{Further rules for TabPFN-3.5.}

\begin{itemize}
\item
  \textbf{Why the standard checkpoint, not Fast (timing).}
  Both TabPFN-3.5 checkpoints were timed on the GPU on 4
  scenarios x 2 replicates of the smoke stream, with the
  run's prediction rows
  (\texttt{results/tabpfn35\_fast\_timing.csv}):

  {\def\LTcaptype{none} % do not increment counter
  \begin{longtable}[]{@{}
    >{\raggedright\arraybackslash}p{(\linewidth - 4\tabcolsep) * \real{0.3333}}
    >{\raggedright\arraybackslash}p{(\linewidth - 4\tabcolsep) * \real{0.3333}}
    >{\raggedright\arraybackslash}p{(\linewidth - 4\tabcolsep) * \real{0.3333}}@{}}
  \toprule\noalign{}
  \begin{minipage}[b]{\linewidth}\raggedright
  Scenario (predicted rows)
  \end{minipage} &
  \begin{minipage}[b]{\linewidth}\raggedright
  Standard, s per fit
  \end{minipage} &
  \begin{minipage}[b]{\linewidth}\raggedright
  Fast, s per fit
  \end{minipage} \\
  \midrule\noalign{}
  \endhead
  \bottomrule\noalign{}
  \endlastfoot
  S1\_pi03\_n200 (4,334) & 23.7-30.6 & 5.2-5.6 \\
  S2\_pi10\_n500 (2,000) & 14.5-14.6 & 3.2-3.3 \\
  S3\_pi30\_n1000 (2,000) & 19.0 & 4.1-4.2 \\
  S4\_pi03\_n3000 (4,334) & 57.4-57.7 & 11.9-12.0 \\
  \end{longtable}
  }

  \begin{itemize}
  \tightlist
  \item
    Fast is 4.5-5.4 times faster. Peak GPU memory was
    2.1-2.9 GB for the standard checkpoint and 1.0-2.2 GB
    for Fast.
  \item
    With the standard checkpoint, schedule P4 (Section 9.3)
    finishes in about 3 days.
  \item
    The plan's budget fallback ``use the Fast variant''
    therefore does not apply.
  \item
    The main arm is the standard checkpoint: the default of
    \texttt{ModelVersion.V3\_5}, and the model the vendor
    ships as TabPFN-3.5.
  \item
    The Fast checkpoint
    (\texttt{tabpfn-v3.5-fast-20260909.safetensors}; SHA-256
    in Section 13) is recorded but not used.
  \end{itemize}
\item
  \textbf{Device.} TabPFN-3.5 runs on the GPU for some
  scenario groups and on CPU workers for others (Section
  9.3). The device is recorded per replicate. On the same
  replicate, CPU and GPU predictions differ by at most
  2.6e-4 (mean 2.9e-5).
\item
  \textbf{The temperature arm's device.} \texttt{tabpfn\_t1}
  runs on the GPU. For its 8 scenarios (pi = 10 \%), the
  main \texttt{tabpfn} arm runs on a CPU worker. The two
  arms therefore differ in temperature (a logit scale of
  1/0.9) and, negligibly, in device.
\end{itemize}

\textbf{Further rules for Mitra-v2.}

\begin{itemize}
\tightlist
\item
  \textbf{Precision.} Mitra-v2 keeps its \textbf{shipped
  bfloat16 autocast}. The GTX 1660 Ti (Turing) has no
  bfloat16 hardware, so the arithmetic is emulated. Two
  support orders change single predictions by 6e-4 on
  average and by at most 0.06. The effect on the true ICI is
  about 1e-4, but E\_max and the tails of OCE are inflated;
  this is reported.
\item
  \textbf{Sign flips.} Mitra's preprocessor flips feature
  signs with the global numpy generator. The wrapper seeds
  that generator from the replicate, so every fit is
  reproducible.
\item
  \textbf{T2, n = 3000.} Mitra-v2 is \textbf{not run at
  T2\_pi10\_n3000}, where it does not fit in 6 GB of GPU
  memory. It is recorded there as ``not run: does not fit in
  6 GB GPU memory''. Its context is never cut. A fit whose
  context AutoGluon cut after an out-of-memory error, in any
  scenario, is recorded as a failure.
\item
  \textbf{Pretraining limits.} The limits of 5,120 support
  rows and 50 features are respected in every scenario.
\end{itemize}

\textbf{Excluded arms and reasons.}

\begin{itemize}
\tightlist
\item
  Jev: a closed hosted decision model with no fitting
  interface and no published weights. It cannot be
  version-pinned, and it would need patient rows sent to a
  third party.
\item
  TabDPT: pretrained on real public tables, a contamination
  risk for the real-data part.
\item
  LimiX-2: does not fit the 6 GB GPU.
\item
  Nori: regression only.
\item
  Zero-shot large language models: shown not to work for
  this task in earlier work.
\end{itemize}

\subsubsection{5.3 External goodness-of-fit battery (R,
CPU)}\label{external-goodness-of-fit-battery-r-cpu}

\textbf{Software.} \texttt{ebrahim.gof} 2.10.0, built from
commit \textbf{19480aa} of the 2.10.0 release branch (C21).

\begin{itemize}
\tightlist
\item
  \textbf{The build is checked, not only the version
  number.} Before the freeze, every one of the 189 functions
  of the installed namespace was compared with the R sources
  of commit 19480aa: 0 differences
  (\texttt{R/verify\_gof\_build.R}). The layer stores the
  SHA-256 of that code and recomputes it from the installed
  package at every start. Any other build is refused,
  including a build of 2c77f3c, which carries the same
  version number.
\item
  Every output row records ``2.10.0@19480aa''.
\end{itemize}

\textbf{What is run on each evaluation set.} For every arm,
replicate and scenario (subject to the schedule of Section
6.3):

\begin{enumerate}
\def\labelenumi{\arabic{enumi}.}
\tightlist
\item
  \textbf{\texttt{run.all.external(y,\ p\_hat,\ G\ =\ 10,\ X)}
  on the observed Y.} Its rows:

  \begin{itemize}
  \tightlist
  \item
    EDGE (G = 10), the primary EDGE row, and EDGE (G =
    auto);
  \item
    the Cox recalibration likelihood-ratio test;
  \item
    calibration-in-the-large;
  \item
    Spiegelhalter z;
  \item
    Osius-Rojek (external);
  \item
    GiViTI (external mode, via \texttt{callr});
  \item
    Hosmer-Lemeshow (external, chi-square with G df);
  \item
    Stukel (offset), the score test (below);
  \item
    le Cessie (external), run when the evaluation covariates
    are stored (replicates 1-R);
  \item
    three descriptive rows.
  \end{itemize}
\item
  \textbf{\texttt{edge.gof(...,\ external\ =\ TRUE)}.} It is
  kept as a cross-check of the battery's EDGE rows.
\item
  \textbf{The same battery on an exact null outcome Y*
  \textasciitilde{} Bernoulli(p-hat).} Y* is drawn from the
  arm's own seed stream and stored with the predictions.
  Under Y*, the arm's predictions are calibrated by
  construction, so the rejection rate is the size of each
  test \emph{on that arm's prediction distribution}.
\item
  \textbf{\texttt{deepgof1.external(y,\ p\_hat,\ X,\ B\ =\ 99)}
  on the observed Y only}, for FLIC and the three TFMs.

  \begin{itemize}
  \tightlist
  \item
    The ``combined'' reading is used, except in designs with
    more than 12 predictors (scenario T2, and the random
    designs with p \textgreater{} 12; Section 3.9), where
    only the \textbf{``axes'' reading} is used. All pairs of
    30 covariates (435 pairs) take about 230 s per call; the
    axes reading takes under 1 s.
  \item
    DeepGOF-1 is an exact Monte Carlo test. Its size was
    checked in the auxiliary size study.
  \end{itemize}
\item
  \textbf{\texttt{localize.external(y,\ p\_hat,\ X,\ M\ =\ 999)}}
  (Section 5.4).
\end{enumerate}

\textbf{The size study (review D5).}

\begin{itemize}
\tightlist
\item
  \textbf{Main size study: the oracle arm.} In every
  scenario, the oracle arm's real Y and its Y* are exact
  nulls, at the main design's DGP. Over replicates 1-500
  this gives 1,000 null evaluation sets per scenario, and
  the full battery is run on all of them.
\item
  \textbf{Each arm's Y*} (replicates 1-R) gives the size on
  that arm's own prediction distribution.
\item
  \textbf{Auxiliary.} The earlier study (2,000 replicates x
  12 scenarios, its own coefficients) is reported as an
  auxiliary check.
\end{itemize}

\textbf{Stukel (offset): the score form, ranked like every
other test (C21).}

\begin{itemize}
\tightlist
\item
  \textbf{The problem.} Up to commit 2c77f3c the row was the
  likelihood-ratio test of Stukel's two terms added to the
  frozen logit. The auxiliary study found it liberal: 0.063
  pooled at N = 3,334, and 0.0745 at S1, pi = 10 \%. A
  separate null simulation (4,000 replicates per cell)
  traced the cause. The upper-tail term, 0.5 lp\^{}2 for lp
  \textgreater= 0, is carried by only the 10-20 records with
  p \textgreater= 0.5 when events are rare, and its
  likelihood-ratio statistic is then not chi-square (size
  0.078 at pi = 10 \% with normal risks).
\item
  \textbf{The change.} From commit 19480aa the row is the
  \textbf{score test of the same two terms at the null}, U'
  I\^{}- U with U = Z'(y - p) and I = Z' diag(p(1 - p)) Z.
  Nothing is fitted. The columns are scaled to their
  information and the inverse is rank-aware. When no
  prediction reaches 0.5 the upper term is empty and the
  test has 1 df; the row's Note says ``score test''.
\item
  \textbf{Its size.} In the null design and with the seeds
  of the earlier Stukel study (4,000 nulls per cell;
  predictions equal to the true risks; pi = 3, 10 and 30 \%
  at N = 3,334, 1,000 and 1,000; normal and t3 risks), the
  package's own row rejects at 0.035-0.053 at the 0.05 level
  (MCSE 0.003-0.0035), with a p-value in all 24,000
  replicates (\texttt{R/stukel\_score\_size.R},
  \texttt{results/stukel\_score\_size.csv}). The author's
  own check gave 0.037-0.055. The lowest value, 0.035 at pi
  = 3 \% with normal risks, is conservative, not liberal:
  there the upper term is empty in 37 \% of the replicates
  (1 df). Under t3 risks some predictions reach 1 in double
  precision; they are first moved to 1 - 1e-9, as the test
  layer does.
\item
  \textbf{Consequence.} Stukel is treated like every other
  test: it is in the size question (Q4a, Section 7.4) with
  no exception, in the detection curves and the ranking of
  tests, and in the summary tables. As for every test, its
  size on the main design comes from the oracle nulls and
  from each arm's Y*, and every table that shows a test's
  power also shows its size.
\end{itemize}

\textbf{Validation-size subsamples.}

\begin{itemize}
\tightlist
\item
  The battery (no DeepGOF-1, no localization) is also run on
  the nested e50 subsample at every event rate, and on the
  e200 subsample at pi = 30 \% only (Section 3.7).
\item
  It is run on the observed Y and on Y*, for every arm
  except the ladder columns.
\item
  This gives detection and size against validation size.
\end{itemize}

\subsubsection{5.4 Localization}\label{localization}

\begin{itemize}
\tightlist
\item
  \textbf{Call.}
  \texttt{ebrahim.gof::localize.external(y,\ p\_hat,\ X,\ M\ =\ 999,\ seed\ =\ \textless{}per-replicate\ seed\textgreater{})},
  with familywise error control at 0.05. It names zero or
  more of the groups INTERCEPT, SLOPE, LINK and COV.
\item
  \textbf{Where it runs.}

  \begin{itemize}
  \tightlist
  \item
    The oracle, FLIC, GAM, the TFMs and the temperature
    arms, in every scenario where the covariates are stored
    (replicates 1-R).
  \item
    The primary analysis is on S3 and T3.
  \item
    The oracle rows (real Y and Y*) give the familywise
    error rate.
  \end{itemize}
\item
  \textbf{T3's factor.} T3's 6-level factor enters as 5
  dummy columns.
\end{itemize}

\subsubsection{5.5 Recalibration
ladder}\label{recalibration-ladder}

For FLIC and the three TFMs, with a local sample of m in
\{100, 500\} rows taken from the recalibration set:

{\def\LTcaptype{none} % do not increment counter
\begin{longtable}[]{@{}
  >{\raggedright\arraybackslash}p{(\linewidth - 2\tabcolsep) * \real{0.2000}}
  >{\raggedright\arraybackslash}p{(\linewidth - 2\tabcolsep) * \real{0.8000}}@{}}
\toprule\noalign{}
\begin{minipage}[b]{\linewidth}\raggedright
Level
\end{minipage} & \begin{minipage}[b]{\linewidth}\raggedright
Method
\end{minipage} \\
\midrule\noalign{}
\endhead
\bottomrule\noalign{}
\endlastfoot
R0 & none \\
R1 & intercept only:
\texttt{glm(y\ \textasciitilde{}\ 1,\ offset\ =\ logit(p\_hat))} \\
R2 & logistic recalibration:
\texttt{glm(y\ \textasciitilde{}\ logit(p\_hat))} \\
R3 & isotonic regression (pool-adjacent-violators), with
predictions clipped as in Section 4 \\
\end{longtable}
}

After each level, the oracle estimands (Sections 4.1-4.2)
and the battery on the observed Y (Section 5.3, item 1) are
computed on the evaluation set. A level whose local sample
has a single outcome class (m = 100 at pi = 3 \%) is
recorded as failed.

\textbf{Test-guided rule (evaluated, not assumed).}

\begin{enumerate}
\def\labelenumi{\arabic{enumi}.}
\tightlist
\item
  If the Cox recalibration test on the R0 predictions
  rejects at 0.05 on the evaluation set, recalibrate with R2
  fitted on the local sample (m = 500).
\item
  If EDGE (G = 10) still rejects at 0.05 on the evaluation
  set with the R2 predictions, flag ``recalibration
  insufficient''.
\end{enumerate}

\begin{center}\rule{0.5\linewidth}{0.5pt}\end{center}

\subsection{6. Performance measures and Monte Carlo
error}\label{performance-measures-and-monte-carlo-error}

\subsubsection{6.1 Per scenario and
arm}\label{per-scenario-and-arm}

\begin{itemize}
\tightlist
\item
  \textbf{Summaries.} The mean, median, inter-quartile
  range, and 10th and 90th percentiles of each estimand
  across replicates.

  \begin{itemize}
  \tightlist
  \item
    Medians and IQRs are the main summaries for a* and b*,
    because single separated fits can produce extreme
    values. In the pilot, ridge's mean b* was 19.7, driven
    by one replicate.
  \end{itemize}
\item
  \textbf{Monte Carlo standard errors.}

  \begin{itemize}
  \tightlist
  \item
    For every mean: SD / sqrt(R).
  \item
    For every proportion: sqrt(p(1 - p) / R).
  \item
    For medians: a distribution-free 95 \% confidence
    interval from order statistics, and a bootstrap MCSE.
  \end{itemize}
\item
  \textbf{Paired differences against FLIC} (each arm minus
  FLIC on the same replicates), with the MCSE of the paired
  difference.
\item
  \textbf{Classes.} The proportion of replicates in each
  calibration class, with MCSE.
\item
  \textbf{Tests.} The rejection rate of every test under the
  observed Y (detection) and under Y* (size), with MCSE.
\item
  \textbf{Fit counts.} Counts of failed, degenerate and
  not-run fits, of \texttt{platt\_fallback}, and of
  \texttt{lambda\_at\_edge} (Section 8).
\end{itemize}

\subsubsection{6.2 Detection curves}\label{detection-curves}

\begin{itemize}
\tightlist
\item
  \textbf{Model.} For each test, P(reject at 0.05) is
  modelled against the true ICI with a logistic GAM
  (\texttt{mgcv}, \texttt{s(true\ ICI,\ k\ =\ 10)}, REML).
\item
  \textbf{Data.} It is fitted across replicates, arms and
  scenarios at N\_eval, with a separate curve for each test.
\item
  \textbf{Uncertainty.} A cluster bootstrap over replicates
  (500 resamples).
\item
  \textbf{Every test of the battery} has its curve and its
  rank, Stukel (offset) included (Section 5.3, C21).
\end{itemize}

\subsubsection{6.3 Replicates (schedule P4) and achieved
MCSE}\label{replicates-schedule-p4-and-achieved-mcse}

{\def\LTcaptype{none} % do not increment counter
\begin{longtable}[]{@{}
  >{\raggedright\arraybackslash}p{(\linewidth - 6\tabcolsep) * \real{0.2500}}
  >{\raggedright\arraybackslash}p{(\linewidth - 6\tabcolsep) * \real{0.1833}}
  >{\raggedright\arraybackslash}p{(\linewidth - 6\tabcolsep) * \real{0.1833}}
  >{\raggedright\arraybackslash}p{(\linewidth - 6\tabcolsep) * \real{0.3833}}@{}}
\toprule\noalign{}
\begin{minipage}[b]{\linewidth}\raggedright
Arms
\end{minipage} & \begin{minipage}[b]{\linewidth}\raggedright
Scenarios
\end{minipage} & \begin{minipage}[b]{\linewidth}\raggedright
R
\end{minipage} & \begin{minipage}[b]{\linewidth}\raggedright
MCSE achieved (pilot SDs)
\end{minipage} \\
\midrule\noalign{}
\endhead
\bottomrule\noalign{}
\endlastfoot
TabPFN-3.5, TabICLv2, Mitra-v2 & all 58 (Mitra-v2 not run at
T2\_pi10\_n3000) & 200 at n \textless= 1000; \textbf{150 at
n = 3000} & mean true ICI \textless= 0.0007 (0.0009 at n =
3000), against a target of 0.002. Rejection rate at 0.05:
0.0154 (0.0178 at n = 3000), against a target of 0.015. At
0.5: 0.035 (0.041). \\
Oracle, LR-MLE, FLIC, Ridge-CV, GAM, LightGBM (+ raw) & all
58 & 500 & mean true ICI \textless= 0.0005 for FLIC and
0.0011 for GAM at n = 200, pi = 3 \% (target 0.001) \\
Real-TabPFN-2 legacy & 16 core scenarios at pi = 10 \% & 100
& descriptive \\
\texttt{tabpfn\_t1}, \texttt{tabicl\_t1} & S1-S4 at pi = 10
\%, n in \{200, 3000\} & 100 & descriptive \\
Oracle size study (oracle real Y + Y*) & all 58 & 500 x 2 =
1,000 nulls & 0.0069 per scenario, 0.0009 pooled \\
Random-design tier: oracle, the classical arms, TabPFN-3.5,
TabICLv2, Mitra-v2 & D = 1,000 random designs (Section 3.9)
& 1 per design & the precision of the meta-model (Section
6.4): a simultaneous 95 \% band of about +/- 0.007 on the
true-ICI difference against FLIC at a typical design
(Section 6.4) \\
\end{longtable}
}

\begin{itemize}
\tightlist
\item
  \textbf{The n = 3000 shortfall.} At n = 3000, R = 150
  misses the rejection-rate target by about 15 \% (0.0178
  against 0.015). This is the plan's first budget fallback,
  adopted before the freeze.
\item
  \textbf{Common replicates.}

  \begin{itemize}
  \tightlist
  \item
    Every paired comparison between a TFM and a classical
    arm uses the replicates that both have (1-200, or 1-150
    at n = 3000).
  \item
    The 500-replicate summaries of the classical arms are
    reported separately.
  \end{itemize}
\item
  \textbf{Battery schedule.}

  \begin{itemize}
  \tightlist
  \item
    The battery runs on every arm for replicates 1-R.
  \item
    It runs on the oracle only (observed Y and Y*) for
    replicates R+1 to 500.
  \item
    No pre-registered estimand needs the battery on the
    classical arms beyond replicate R; their oracle
    estimands come from all 500.
  \end{itemize}
\end{itemize}

\textbf{Reporting rule.} No statement of ``no difference''
or ``holds size'' is written without first giving the MCSE
(or confidence interval) of the quantity. A scenario where
the quantity is at a floor or a ceiling is reported as
uninformative for that comparison.

\subsubsection{6.4 Meta-model of the random-design
tier}\label{meta-model-of-the-random-design-tier}

The random designs (Section 3.9) are analysed with a
pre-specified meta-model, in R \texttt{mgcv} (REML).

\textbf{Responses.} For each design d and each TFM m
(TabPFN-3.5, TabICLv2, Mitra-v2), paired with FLIC on the
same data:

\begin{itemize}
\tightlist
\item
  D\_ICI(d, m) = true ICI(m) - true ICI(FLIC);
\item
  D\_logb(d, m) = \textbar log b*(m)\textbar{} -
  \textbar log b*(FLIC)\textbar, where \textbar log
  b*\textbar{} is capped at log 10 (b* \textless= 0 or no
  finite b* counts as log 10). The meta-model uses the
  values themselves, so it needs this finite cap; the
  order-based summaries of Section 7.1 place such replicates
  above every finite value instead.
\end{itemize}

\textbf{Model.} One GAM per response, fitted to all designs
where both arms produced usable predictions (Section 8.2):

\begin{verbatim}
D ~ model + s(log n, by = model, k = 6) + s(log pi, by = model, k = 6)
      + s(p, by = model, k = 5) + s(AUC target, by = model, k = 5)
      + s(w_nl, by = model, k = 4) + s(w_int, by = model, k = 4)
      + s(w_link, by = model, k = 4) + model:rho + model:discrete fraction
      + s(design, bs = "re")
\end{verbatim}

\begin{itemize}
\tightlist
\item
  \texttt{model} is the TFM factor; each TFM gets its own
  partial effects.
\item
  The design random effect pairs the three TFMs that share a
  design's data.
\item
  The family is the scaled t (\texttt{scat}), so that a few
  separated FLIC fits cannot dominate.
\end{itemize}

\textbf{Outputs.}

\begin{itemize}
\tightlist
\item
  Partial-effect plots of every smooth, per TFM, with 95 \%
  \textbf{simultaneous} bands. The bands come from 10,000
  draws of the coefficients from their approximate
  posterior, with the critical value of the maximum
  standardized deviation over the plotted grid.
\item
  \textbf{Regions.} For each TFM, every random design's own
  setting is classified as ``TFM better'' (the simultaneous
  band of the predicted D\_ICI lies below -0.01), ``TFM
  worse'' (above +0.01) or ``within 0.01''. The shares are
  reported, with maps on (log n, log pi) at the median of
  the other factors, and on the two factors with the largest
  partial effects.
\item
  The same maps for D\_logb, with a margin of log 1.1.
\item
  \textbf{Detection.} For every test and every arm, the
  rejection at 0.05 on the observed Y (0/1) is modelled with
  the same terms and a logit link (binomial family, no
  design random effect, one model per arm and test). The
  size on the oracle and on each arm's Y* is modelled the
  same way, and every region where the simultaneous band of
  the predicted size excludes {[}0.04, 0.06{]} is reported.
\item
  \textbf{Mitra-v2} is modelled only where it ran (Section
  3.9); the region where it did not fit is shown on every
  map.
\end{itemize}

\textbf{Precision.} One replicate per design is enough
because the meta-model pools neighbouring designs. Take a
residual SD of 0.012 for D\_ICI (the pilot's largest SD of a
paired true-ICI difference) and about 40 effective degrees
of freedom per TFM. With D = 1,000, the pointwise SE of the
predicted difference at a typical design is then about 0.012
x sqrt(40 / 1,000) = 0.0024, and the simultaneous band is
about three times wider (+/- 0.007). A difference larger
than about 0.017 is therefore classified as ``better'' or
``worse'' at the 0.01 margin. At the edges of the design
space the bands are wider, and the maps show them.

\begin{center}\rule{0.5\linewidth}{0.5pt}\end{center}

\subsection{7. Research questions and reporting
rules}\label{research-questions-and-reporting-rules}

This is an estimation-focused protocol. It states no
confirmatory hypothesis and predicts the direction of no
result. Each research question below fixes, before any
main-run result is seen, the estimands, the comparison, the
summary with its Monte Carlo standard error (MCSE), and the
reporting rule that decides when a model is called better,
worse or equivalent. The results then determine the
conclusions, and the rules cannot change after the data are
seen.

\subsubsection{7.0 Rules common to every
question}\label{rules-common-to-every-question}

\begin{itemize}
\tightlist
\item
  \textbf{``TFM''.} The word means each of the three TFM
  arms judged separately: TabPFN-3.5 (\texttt{tabpfn}),
  Mitra-v2 and TabICLv2 at their shipped settings. The
  legacy and temperature arms are not part of the main
  comparison set; they are reported beside it.
\item
  \textbf{Comparator.} Every comparison is paired against
  FLIC: the TFM minus FLIC on the same replicates. It uses
  the replicates where both arms produced usable predictions
  (Section 8.2).
\item
  \textbf{Summaries.}

  \begin{itemize}
  \tightlist
  \item
    For a difference in true ICI: the mean paired
    difference, with MCSE = SD / sqrt(R) and the 95 \%
    interval mean +/- 1.96 MCSE.
  \item
    For differences in \textbar log b*\textbar{} and
    \textbar a*\textbar, which single separated fits can
    make extreme: the Hodges-Lehmann (HL) estimate of the
    paired differences, with its distribution-free 95 \%
    interval (from the signed-rank statistic) and a
    bootstrap MCSE.
  \end{itemize}
\item
  \textbf{Classification.} For a paired difference where a
  negative value favours the TFM, and a margin delta:

  \begin{itemize}
  \tightlist
  \item
    \textbf{better} than FLIC: the whole 95 \% interval lies
    below -delta;
  \item
    \textbf{worse} than FLIC: the whole interval lies above
    +delta;
  \item
    \textbf{equivalent} to FLIC: the whole interval lies
    inside (-delta, +delta);
  \item
    \textbf{inconclusive}: otherwise.
  \end{itemize}
\item
  \textbf{Margins}, fixed now:

  \begin{itemize}
  \tightlist
  \item
    true ICI: delta = 0.01, the margin of Section 6.4; at pi
    = 3 \% it is multiplied by 0.3 (delta = 0.003), by the
    rule for probability-scale thresholds of Section 4.5;
  \item
    \textbar log b*\textbar: delta = log 1.1, the slope
    margin of Section 6.4;
  \item
    \textbar a*\textbar: delta = 0.05, half the
    weak-calibration tolerance on a* (Section 4.5).
  \end{itemize}
\item
  \textbf{MCSE requirement.} The reporting rule of Section
  6.3 applies to every question: no ``equivalent'', ``no
  difference'' or ``holds its size'' is written without the
  interval or MCSE of the quantity, and a cell at a floor or
  a ceiling is reported as uninformative for that
  comparison.
\item
  \textbf{No p-value decides a label.} The intervals are 95
  \% intervals, not adjusted for multiplicity. Each question
  reports how many cells fall in each class, so that the
  reader can weigh the number of cells.
\item
  \textbf{Reporting.} Every label is reported in a
  summary-of-findings table with one row per question,
  whatever it is.
\end{itemize}

\subsubsection{7.1 Q1. How far from 1, and in which
direction, is each TFM's calibration slope at small
n?}\label{q1.-how-far-from-1-and-in-which-direction-is-each-tfms-calibration-slope-at-small-n}

\begin{itemize}
\tightlist
\item
  \textbf{Estimands.} The true slope b* of each TFM and of
  FLIC (Section 4.1), and the paired difference d =
  \textbar log b*(TFM)\textbar{} - \textbar log
  b*(FLIC)\textbar. A replicate with b* \textless= 0 or no
  finite b* gets \textbar log b*\textbar{} = 10, which
  places it above every finite value.
\item
  \textbf{Main cells.} pi = 10 \% and n in \{200, 500\}, in
  the four shapes: 8 cells per TFM. The same summaries are
  reported in all 48 core scenarios.
\item
  \textbf{Summaries.} The HL estimate of d with its 95 \%
  interval and bootstrap MCSE; the median b* of each arm
  with its distribution-free 95 \% interval (Section 6.1).
\item
  \textbf{Reporting rules.}

  \begin{itemize}
  \tightlist
  \item
    \emph{Closeness to 1.} In each cell, the TFM is better
    than, worse than or equivalent to FLIC, or inconclusive,
    by Section 7.0 with delta = log 1.1.
  \item
    \emph{Direction.} In each cell, an arm is called
    \textbf{under-confident} when the 95 \% interval of its
    median b* lies above 1.1, \textbf{over-confident} when
    it lies below 0.9, \textbf{close to 1} when it lies
    inside {[}0.9, 1.1{]} (the slope tolerance of weak
    calibration, Section 4.5), and \textbf{unclear}
    otherwise. FLIC is labelled in the same way.
  \item
    For each TFM, the number of the 8 main cells in each
    class is reported.
  \end{itemize}
\end{itemize}

\subsubsection{7.2 Q2. At what development size does each
TFM reach calibration, compared with
FLIC?}\label{q2.-at-what-development-size-does-each-tfm-reach-calibration-compared-with-flic}

\begin{itemize}
\tightlist
\item
  \textbf{Estimands.} The minimum n for calibration (Section
  4.5) of each arm, in each shape, at each event rate and
  for each class; and the proportion of replicates that meet
  each class at each n, with MCSE.
\item
  \textbf{Comparison.} For each TFM, shape and event rate:
  the TFM's minimum n for weak calibration against FLIC's,
  both from the common replicates. ``\textgreater{} 3000''
  ranks above 3000.
\item
  \textbf{Reporting rules.}

  \begin{itemize}
  \tightlist
  \item
    The TFM \textbf{needs a smaller n} than FLIC,
    \textbf{the same n} or \textbf{a larger n}, by comparing
    the two minimum n.
  \item
    The comparison is repeated with the 80 \% criterion
    replaced by its lower 95 \% limit (proportion - 1.96
    MCSE \textgreater= 0.80). An ordering is called
    \textbf{clear} only when both versions give it;
    otherwise it is called \textbf{sensitive to Monte Carlo
    error}.
  \item
    The comparison is also repeated at half and at double
    every threshold, and reported beside the main one.
  \item
    Results are reported for all four shapes and all three
    event rates; no shape is singled out in advance.
  \end{itemize}
\end{itemize}

\subsubsection{7.3 Q3. How large is each TFM's
calibration-in-the-large error when events are
rare?}\label{q3.-how-large-is-each-tfms-calibration-in-the-large-error-when-events-are-rare}

\begin{itemize}
\tightlist
\item
  \textbf{Estimands.} The true intercept a* (Section 4.1),
  the paired difference d = \textbar a*\textbar(TFM) -
  \textbar a*\textbar(FLIC), and mean p-hat minus mean p*
  (Section 4.2).
\item
  \textbf{Cells.} pi = 3 \%, the four shapes and all four n:
  16 cells per TFM. The cells at n \textgreater= 1000 are
  shown first, because the pilot suggested that the shipped
  softmax temperature matters most there (Section 0.2, C2).
  The cells at n \textless= 500 are reported with the same
  rule. The same summaries are also reported at pi = 10 \%
  and 30 \%.
\item
  \textbf{Summaries.} The HL estimate of d with its 95 \%
  interval and bootstrap MCSE; the median a* of each arm
  with its 95 \% interval.
\item
  \textbf{Reporting rule.} Better, worse, equivalent or
  inconclusive by Section 7.0, with delta = 0.05. Each TFM's
  shipped softmax temperature is shown beside its result.
\item
  \textbf{Temperature arms.} At pi = 10 \%, the paired
  effect of the temperature (\texttt{tabpfn\_t1} minus
  \texttt{tabpfn}, and \texttt{tabicl\_t1} minus
  \texttt{tabicl}) on a*, b*, mean risk and true ICI is
  reported with the same summaries and margins, as a direct
  estimate of the mechanism.
\end{itemize}

\subsubsection{7.4 Q4. How do the goodness-of-fit tests
behave on TFM
predictions?}\label{q4.-how-do-the-goodness-of-fit-tests-behave-on-tfm-predictions}

\emph{Q4a. Size.}

\begin{itemize}
\tightlist
\item
  \textbf{Estimand.} The rejection rate at 0.05 of every
  battery test under the exact null Y* \textasciitilde{}
  Bernoulli(p-hat), for every arm, with MCSE. For the
  oracle, its real Y and Y* over replicates 1-500 are used
  (Section 5.3). Stukel (offset), now the score test, is
  included with no exception (C21).
\item
  \textbf{Reporting rule}, for each (test, arm):

  \begin{itemize}
  \tightlist
  \item
    \textbf{holds its size}: the pooled size over all
    scenarios where the arm was run lies in {[}0.04,
    0.06{]}, and the per-scenario size lies inside its 99 \%
    Monte Carlo band (0.05 +/- 2.576 MCSE) in at least 95 \%
    of scenarios (about 1 \% of scenarios are expected
    outside the band by chance);
  \item
    \textbf{liberal}: the pooled size is above 0.06;
  \item
    \textbf{conservative}: the pooled size is below 0.04;
  \item
    \textbf{uneven}: the pooled size lies in {[}0.04,
    0.06{]}, but fewer than 95 \% of scenarios are inside
    their band.
  \end{itemize}
\item
  Every (test, arm, scenario) outside its band is listed.
\end{itemize}

\emph{Q4b. Detection.}

\begin{itemize}
\tightlist
\item
  \textbf{Estimands.} For every test, the detection curve
  P(reject at 0.05 \textbar{} true ICI) of Section 6.2, and
  the rank of the tests.
\item
  \textbf{Main comparison.} EDGE (G = 10) against
  Hosmer-Lemeshow (external): the difference P(reject
  \textbar{} EDGE) - P(reject \textbar{} HL) at true ICI in
  \{0.01, 0.02, 0.03, 0.05\}, from curves fitted on the same
  replicates, with cluster-bootstrap 95 \% intervals.
\item
  \textbf{Reporting rule}, at each of the four points, with
  a margin of 2 percentage points: EDGE \textbf{detects more
  often} when the interval lies above +0.02, \textbf{less
  often} when it lies below -0.02, \textbf{equivalently}
  when it lies inside (-0.02, +0.02), and the comparison is
  \textbf{inconclusive} otherwise.
\item
  Every table and figure that shows a test's detection also
  shows its size (Section 5.3).
\end{itemize}

\subsubsection{7.5 Q5. How well does recalibration repair
each TFM, and does the test-guided rule pick the right
cases?}\label{q5.-how-well-does-recalibration-repair-each-tfm-and-does-the-test-guided-rule-pick-the-right-cases}

\emph{Q5a. Repair.}

\begin{itemize}
\tightlist
\item
  \textbf{Estimands.} The mean true ICI of each arm after
  each level of the ladder (Section 5.5), with MCSE.
\item
  \textbf{Comparisons.}

  \begin{itemize}
  \tightlist
  \item
    \emph{Against FLIC.} The TFM after R2 (m = 500) minus
    FLIC after R2 (m = 500), paired, labelled by Section 7.0
    with the true-ICI margin.
  \item
    \emph{Against the best arm.} In each core scenario, the
    best arm is the one with the smallest mean true ICI
    among FLIC and the three TFMs after R2 (m = 500), and
    LR-MLE, Ridge, GAM and LightGBM (Platt) as fitted. The
    statistic is the TFM's mean true ICI after R2 minus the
    best arm's mean true ICI, with the paired MCSE.
  \end{itemize}
\item
  \textbf{Reporting rule against the best arm.} In a
  scenario, the TFM is \textbf{within 0.005 of the best arm}
  when the upper 95 \% limit of the difference is \textless=
  0.005, \textbf{more than 0.005 above it} when the lower
  limit is \textgreater{} 0.005, and \textbf{unclear}
  otherwise. For each TFM, the counts over the 48 core
  scenarios are reported.
\end{itemize}

\emph{Q5b. The test-guided rule.}

\begin{itemize}
\tightlist
\item
  \textbf{Truth.} A replicate is ``truly insufficient'' if
  its true ICI after R2 (m = 500) is \textgreater{} 0.015,
  or \textgreater{} 0.0045 at pi = 3 \% (Section 4.5).
\item
  \textbf{Estimands.} PPV = P(truly insufficient \textbar{}
  flagged) and NPV = P(not truly insufficient \textbar{} Cox
  rejected and not flagged), pooled over the core scenarios,
  for each TFM and for FLIC, with cluster-bootstrap 95 \%
  intervals.
\item
  \textbf{Reporting rule.} The PPV is called \textbf{at
  least 0.8} when its lower 95 \% limit is \textgreater=
  0.8, \textbf{below 0.8} when its upper limit is
  \textless{} 0.8, and \textbf{unclear} otherwise. If fewer
  than 50 replicates are flagged for an arm, its PPV is
  reported as ``not assessable'' (its MCSE would exceed
  0.057). The NPV is reported with its interval and no
  label.
\end{itemize}

\subsubsection{7.6 Q6. Do the grid findings carry over to
randomly drawn designs? (robustness
check)}\label{q6.-do-the-grid-findings-carry-over-to-randomly-drawn-designs-robustness-check}

Q6 is answered by a pre-specified robustness check. It
checks whether the core design's results on the true ICI
transfer to the random-design tier (Section 3.9). It is a
check of transfer, not a hypothesis.

\begin{itemize}
\tightlist
\item
  \textbf{Transfer check.} For each TFM and each of the 48
  core scenarios, the meta-model of Section 6.4 predicts
  D\_ICI at the scenario's own coordinates: its n and pi, p
  = 8, AUC target = its oracle AUC, rho = 0.018 (the mean
  pairwise correlation of the core covariates: one pair at
  0.5 among 28), discrete fraction 0.25, and the severities
  of its shape (S1: all 0; S2: w\_nl = 0.5; S3: w\_int = 1;
  S4: w\_link = 1).
\item
  A scenario \textbf{agrees} when the core design's mean
  paired difference (TFM - FLIC, replicates 1-R) and the
  prediction differ by less than 1.96 sqrt(SE\_meta\^{}2 +
  MCSE\_core\^{}2).
\item
  \textbf{Reporting rule}, per TFM:

  \begin{itemize}
  \tightlist
  \item
    \textbf{the grid transfers}: at least 40 of the 48 core
    scenarios agree (83 \%);
  \item
    \textbf{the grid transfers in part}: 30 to 39 of the 48
    agree;
  \item
    \textbf{the grid does not transfer}: fewer than 30
    agree.
  \end{itemize}
\item
  Mitra-v2 is checked only on the core scenarios inside the
  region where it ran in the random tier.
\item
  \textbf{Reported with the check, whatever it shows:} the
  robustness map, that is, the regions of Section 6.4 where
  each TFM is better or worse than FLIC by more than 0.01;
  and every core scenario that does not agree.
\item
  \textbf{Interpretation, fixed in advance.} If the grid
  does not transfer for a TFM, the core design's fixed
  levels are not typical of the wider design space for that
  TFM. The paper then reports the random-tier maps as the
  general result for that TFM, and the core design as a set
  of special cases.
\end{itemize}

\subsubsection{7.7 Equal prominence for every
outcome}\label{equal-prominence-for-every-outcome}

\begin{itemize}
\tightlist
\item
  If every TFM is at least weakly calibrated in at least 80
  \% of replicates at n = 200, in every shape and at every
  event rate, the paper reports that TFMs remove the
  small-sample calibration problem, with the same prominence
  as the opposite result.
\item
  Whatever direction Q1 finds (under-confident,
  over-confident or close to 1) is reported as the finding,
  with the same prominence.
\end{itemize}

\begin{center}\rule{0.5\linewidth}{0.5pt}\end{center}

\subsection{8. Exclusions, failures and
stopping}\label{exclusions-failures-and-stopping}

\subsubsection{8.1 Stopping}\label{stopping}

\begin{itemize}
\tightlist
\item
  The replicate counts of Section 6.3 are fixed. They are
  not changed after any result of the main run is seen.
\item
  No replicate is redrawn for any reason, including a
  development set with no events.
\end{itemize}

\subsubsection{8.2 Failures, degenerate fits and
pairing}\label{failures-degenerate-fits-and-pairing}

\begin{itemize}
\tightlist
\item
  \textbf{Failure.}

  \begin{itemize}
  \tightlist
  \item
    An arm counts as failed when it raises an error or
    returns a missing, non-finite or out-of-range
    prediction.
  \item
    A single-class development set is recorded as the same
    failure for every arm except the oracle.
  \item
    Failures are logged with the error message and kept.
  \item
    The R arms have a limit of 1,800 s per call. The Python
    arms have no time limit (the pilot's longest fit was 134
    s). A crashed or killed process is restarted, and the
    restart refits only the missing arms (resume).
  \end{itemize}
\item
  \textbf{Not run.} Mitra-v2 at T2\_pi10\_n3000 (Section
  5.2), and in the random designs above its memory limit
  (Section 3.9). It is shown in every table as ``not run''.
\item
  \textbf{Degenerate fit.} A fit that ran is labelled
  \textbf{``degenerate''} (not ``ok'') in either of two
  cases:

  \begin{itemize}
  \tightlist
  \item
    more than 100 of its evaluation predictions lie within
    1e-6 of 0 or 1;
  \item
    LR-MLE shows separation (non-convergence, fitted
    probabilities numerically 0 or 1, or
    \textbar coefficient\textbar{} \textgreater{} 15).
  \end{itemize}

  Degenerate fits are the arm's real output. They are
  scored, tested and kept in every comparison, and counted
  in every table. LR-MLE summaries are also reported without
  them.
\item
  \textbf{Other flags,} counted per arm and scenario:

  \begin{itemize}
  \tightlist
  \item
    \texttt{platt\_fallback} (LightGBM);
  \item
    \texttt{lambda\_at\_edge} (ridge);
  \item
    \texttt{small\_n\_flag} (TabICLv2 at n \textless{} 300);
  \item
    the number of predictions at the boundary.
  \end{itemize}
\item
  \textbf{Pairing rule} (review D8).

  \begin{itemize}
  \tightlist
  \item
    Every comparison between arms uses only the replicates
    where both arms produced usable (ok or degenerate)
    predictions.
  \item
    The known few-event failures are part of the result, not
    hidden:

    \begin{itemize}
    \tightlist
    \item
      ridge fails when the development set has 2 or fewer
      events (glmnet needs 2 events in every CV fold), about
      6 \% of replicates at n = 200, pi = 3 \%;
    \item
      LightGBM fails at 1 or fewer events;
    \item
      LR, FLIC, GAM and the TFMs run with 1 event;
    \item
      the ladder's R1-R3 at m = 100 fail when the 100-row
      local sample has no event;
    \item
      EDGE is ``not run'' on isotonic (R3, m = 100)
      predictions with too few distinct values; such rows
      count neither as rejections nor as non-rejections, and
      the test-guided rule uses R2 only.
    \end{itemize}
  \item
    \textbf{The 5 \% flag.} If an arm fails in more than 5
    \% of the replicates of a scenario, that scenario is
    flagged in all tables, and its results for that arm are
    reported but not used in the labels of Section 7.
    Because failures are not redrawn, this flag is part of
    the result.
  \item
    \textbf{Random designs.} Each design has one replicate,
    so the 5 \% flag does not apply. A failed or not-run arm
    leaves that design out of that arm's meta-model, and the
    failures are mapped over the design space like any other
    outcome.
  \end{itemize}
\end{itemize}

\begin{center}\rule{0.5\linewidth}{0.5pt}\end{center}

\subsection{9. Pilot, smoke test and compute
plan}\label{pilot-smoke-test-and-compute-plan}

\subsubsection{9.1 Phase-1 pilot (before the freeze; smoke
seed
stream)}\label{phase-1-pilot-before-the-freeze-smoke-seed-stream}

See Section 0. The pilot report records:

\begin{itemize}
\tightlist
\item
  times;
\item
  GPU memory;
\item
  the failure rates;
\item
  the calibration summaries;
\item
  the test-layer cost;
\item
  the reproducibility checks: predictions are bit-identical
  across processes and under CPU load, and TabICL differs by
  at most 8e-7.
\end{itemize}

\subsubsection{9.2 Final smoke test (2026-10-07, before the
freeze)}\label{final-smoke-test-2026-10-07-before-the-freeze}

\textbf{What was run.} Every arm on S1\_pi03\_n200 and
T2\_pi10\_n500, 2 replicates each, on the smoke stream. The
files were laid out exactly as in the main run: classical
arms, the GPU process, and TabPFN-3.5 on a CPU worker for
T2. Then the scoring and the test layer: the oracle-only
mode first, then every arm.

\textbf{Results.}

\begin{itemize}
\tightlist
\item
  \textbf{Fits.} All 52 arm fits ran: 49 ``ok'' and 3
  ``degenerate'' (LR-MLE and GAM at S1, GAM at T2), with 0
  failures.
\item
  \textbf{Temperatures and checkpoints.} The recorded
  temperatures were 1.0 (\texttt{tabpfn}), 0.9
  (\texttt{tabpfn\_t1}, \texttt{tabicl},
  \texttt{tabpfn\_legacy}) and 1.0 (\texttt{tabicl\_t1}).
  Both TabPFN-3.5 arms carried the pinned SHA-256.
\item
  \textbf{Test layer.} It wrote gzip files with 0 errors.
  Every arm had both real-Y and Y* rows, including an
  ``Osius-Rojek (external)'' row with a finite p-value (104
  of 104 arm x outcome pairs). DeepGOF-1 used the axes
  reading in T2 (0.2-0.3 s per call).
\item
  \textbf{Pass criteria.} The pass criteria were met: no
  unlogged error, finite predictions, fingerprints asserted,
  and files written and resumable.
\item
  \textbf{Unit tests.} 41 Python tests (in both
  environments), 11 R tests and the kill-and-resume check
  passed.
\end{itemize}

\textbf{Second smoke run (2026-10-07, after C21-C23).} The
package pin moved to commit 19480aa and the random-design
tier was added, so the smoke was repeated on the smoke
stream (\texttt{code/smoke\_freeze\_b.sh},
\texttt{results/smoke\_freeze\_b/SUMMARY.md}):

\begin{itemize}
\tightlist
\item
  \textbf{What was run.} Every arm on S1\_pi03\_n200 and
  T2\_pi10\_n500 (2 replicates each), and the tier's arms on
  the random designs R0001-R0005 (one replicate each), laid
  out as in the main run. TabPFN-3.5 ran on the GPU for four
  designs and on a CPU worker for R0002 ((n + N\_eval) x p =
  128,910). Then the scoring and the test layer: the
  oracle-only mode first, then every arm.
\item
  \textbf{Fits.} Of the 102 fits, 94 were ``ok'', 7
  ``degenerate'' (LR-MLE and GAM in S1 and in two random
  designs, GAM in T2) and 1 ``not run'': Mitra-v2 in R0002
  (n\^{}2 x p = 9.8e7), recorded as ``does not fit in 6 GB
  GPU memory''. There were 0 failures.
\item
  \textbf{Test layer.} 9 gzip files with 0 open errors,
  every one with the reference 2.10.0@19480aa. Every arm had
  real-Y and Y* rows. All 202 Stukel (offset) rows say
  ``score test'' and have a finite p-value (df 2, or 1 where
  no prediction reached 0.5), and all 202 arm x outcome
  pairs have an ``Osius-Rojek (external)'' row with a finite
  p-value. DeepGOF-1 used the axes reading in T2 and in the
  three random designs with p \textgreater{} 12, and the
  combined reading in R0001 (p = 12) and R0005 (p = 8). The
  52 rows marked ``not run'' or ``failed'' all belong to the
  degenerate GAM fits and to the isotonic ladder column
  \texttt{flic\_\_R3\_m100}, as Section 8.2 foresees.
\item
  \textbf{Unit tests.} 48 Python tests (in both
  environments), 16 R tests, the kill-and-resume check,
  \texttt{dgp.py\ check} (the core table is unchanged,
  difference 0), the core-fingerprint test of the family-R
  code and the build check of commit 19480aa passed.
\end{itemize}

\subsubsection{9.3 Schedule P4: streams, devices and
budget}\label{schedule-p4-streams-devices-and-budget}

\textbf{Streams.} Four streams run in parallel on a 12-core
/ 24-thread machine with one GTX 1660 Ti 6 GB. About 5 cores
stay free. A fifth stream, the random-design tier, takes
over their slots as they finish.

{\def\LTcaptype{none} % do not increment counter
\begin{longtable}[]{@{}
  >{\raggedright\arraybackslash}p{(\linewidth - 4\tabcolsep) * \real{0.2333}}
  >{\raggedright\arraybackslash}p{(\linewidth - 4\tabcolsep) * \real{0.1667}}
  >{\raggedright\arraybackslash}p{(\linewidth - 4\tabcolsep) * \real{0.6000}}@{}}
\toprule\noalign{}
\begin{minipage}[b]{\linewidth}\raggedright
Stream
\end{minipage} & \begin{minipage}[b]{\linewidth}\raggedright
Cores
\end{minipage} & \begin{minipage}[b]{\linewidth}\raggedright
Work, in order
\end{minipage} \\
\midrule\noalign{}
\endhead
\bottomrule\noalign{}
\endlastfoot
GPU process (\texttt{-\/-device\ cuda}, 6 threads) & GPU + 1
& TabICLv2 and Mitra-v2 on all scenarios; then TabPFN-3.5 on
the 16 core scenarios at pi = 3 \% and on T2\_pi10\_n500;
then legacy V2; then the two temperature arms \\
CPU worker A (TabPFN-3.5, \texttt{-\/-device\ cpu}, 6
threads) & 6 & T2\_pi10\_n3000 first, then the 16 core
scenarios at pi = 30 \%, then T1 and T3 (4 scenarios) \\
CPU worker B (TabPFN-3.5, \texttt{-\/-device\ cpu}, 6
threads) & 6 & the 16 core scenarios at pi = 10 \%, then T4a
and T4b (4 scenarios) \\
Classical and tests (6 workers) & 6 & the classical arms (R
= 500) first; then the test layer and the scoring, streaming
every pass until the model streams finish; then the
validation-size subsamples \\
Random-design tier (Section 3.9) & the slots it takes over &
when CPU worker B is done: the classical arms of the random
designs (6 workers), then TabPFN-3.5 on half of the CPU
designs; when worker A is done: the other half; when the GPU
process is done: TabICLv2 and Mitra-v2 on every design and
TabPFN-3.5 on the GPU designs; when the classical stream is
done: the test layer and the scoring \\
\end{longtable}
}

\textbf{Device rules.}

\begin{itemize}
\tightlist
\item
  Each scenario's TabPFN-3.5 fits run on one device only.
  The three groups above cover the 58 scenarios exactly once
  (\texttt{code/main\_groups.py} asserts it).
\item
  After about 24 hours the balance is checked:

  \begin{itemize}
  \tightlist
  \item
    a CPU worker that is clearly behind hands its
    \emph{unstarted} scenarios, whole, to the GPU process
    (it is stopped first, so no replicate is fitted twice);
  \item
    a worker that is ahead takes unstarted scenarios from
    the other.
  \end{itemize}
\item
  The device actually used for every scenario is reported in
  the paper. The scenarios fitted on the CPU are listed in
  \texttt{DEVIATIONS.md} only if they differ from this plan.
\item
  The random designs' TabPFN-3.5 device follows the size
  rule of Section 3.9. Those designs are not moved in the
  balance check.
\end{itemize}

\textbf{Budget} (pilot measurements, about 118 GPU-hours
plus about 490 core-hours before balancing):

\begin{itemize}
\tightlist
\item
  about 79 hours for the GPU stream, including about 8
  GPU-hours for the two temperature arms;
\item
  63-77 hours for CPU worker A and 57-71 hours for CPU
  worker B (CPU/GPU time ratio 2.4-3.0);
\item
  the classical arms take about 9 hours of the fourth
  stream;
\item
  the 58 scenarios end after about 3.3-3.4 days (about 82
  hours, with the last pass of the test layer).
\end{itemize}

The DeepGOF-1 axes reading in T2 removes about 70 core-hours
from the pilot's estimate of the test layer (at B = 99).

\textbf{The random-design tier} (D = 1,000). The costs are
measured per fit (\texttt{code/random\_budget.py}, whose
output is \texttt{results/random\_budget.md}). Times such as
``79 h'' are hours from the start of the run.

\begin{itemize}
\tightlist
\item
  \textbf{GPU:} about 6.8 hours, from about 79 h: TabICLv2
  and Mitra-v2 on every design where they run, and
  TabPFN-3.5 on the 821 GPU designs.
\item
  \textbf{TabPFN-3.5 on the 179 CPU designs:} about 5.3
  hours on one 6-thread worker, so about 2.7 hours on each
  of the two workers; done by about 73 h.
\item
  \textbf{The classical arms:} about 4.8 core-hours (about
  50 minutes on 6 workers), from about 64 h.
\item
  \textbf{The test layer and the scoring:} about 13
  core-hours (about 46 core-seconds per design at B = 99).
\item
  \textbf{The basis:} per-fit time models fitted to the
  pilot's and the memory probe's GPU timings (log-scale RMSE
  0.14-0.34), a CPU/GPU time ratio of 3.0 for TabPFN-3.5 on
  6 threads, and the 50-design cost sample for the classical
  arms and the test layer, measured on a loaded machine.
\item
  \textbf{Choice of D.} The tier was to add at most about
  one day, with D \textgreater= 600 and a target of 1,000.
  With D = 1,000 it adds about 4.3 hours (about 11 hours if
  every cost were twice the measured one), so D = 1,000.
\item
  \textbf{New total.} The expected wall time of the whole
  run is about \textbf{3.6 days} (about 86 hours).
\end{itemize}

\textbf{Budget fallbacks of the plan.}

\begin{enumerate}
\def\labelenumi{\arabic{enumi}.}
\tightlist
\item
  R = 150 at n = 3000: \textbf{adopted}.
\item
  the faster variant of TabPFN-3.5: \textbf{not needed}
  (Section 5.2).
\item
  drop the legacy arm: \textbf{not adopted}; the legacy arm
  is kept at R = 100.
\end{enumerate}

\subsubsection{9.4 Engineering}\label{engineering}

\begin{itemize}
\tightlist
\item
  \textbf{Writing replicates.}

  \begin{itemize}
  \tightlist
  \item
    Every replicate is written to disk as soon as it is
    done. The file holds the fingerprints, p*, Y, Y* and the
    p-hat of every arm as float32, and it is written
    atomically.
  \item
    A lock file per replicate makes the streams safe to run
    together.
  \item
    A killed stream resumes with the same command.
  \end{itemize}
\item
  \textbf{File hand-off.} The statistical arms and the test
  layer run in R on the saved files; the TFMs run in Python;
  the hand-off is by file only. The test layer writes one
  gzip-compressed csv per replicate and reads csv or csv.gz.
\item
  \textbf{The start command checks:}

  \begin{itemize}
  \tightlist
  \item
    the software versions of Section 13;
  \item
    the SHA-256 of every checkpoint (it writes
    \texttt{results/environment.json});
  \item
    that the installed \texttt{ebrahim.gof} is the build of
    commit 19480aa (\texttt{R/verify\_gof\_build.R});
  \item
    that the random designs re-solve to the stored table
    (\texttt{code/dgp\_random.py\ check});
  \item
    the test-layer unit tests.
  \end{itemize}
\item
  \textbf{The deposit guard.} The run refuses to start until
  the registration link and timestamp are recorded in
  \texttt{prereg/DEPOSITED.txt}.
\end{itemize}

\begin{center}\rule{0.5\linewidth}{0.5pt}\end{center}

\subsection{10. Real-data analysis
plan}\label{real-data-analysis-plan}

\subsubsection{10.1 Data sets, outcomes and
splits}\label{data-sets-outcomes-and-splits}

Event rates are computed from the files and reported; they
are not taken from the literature. The predictor lists below
are fixed from the codebooks before any outcome is
tabulated, and they are limited to variables available at
the time of prediction.

{\def\LTcaptype{none} % do not increment counter
\begin{longtable}[]{@{}
  >{\raggedright\arraybackslash}p{(\linewidth - 6\tabcolsep) * \real{0.2295}}
  >{\raggedright\arraybackslash}p{(\linewidth - 6\tabcolsep) * \real{0.1803}}
  >{\raggedright\arraybackslash}p{(\linewidth - 6\tabcolsep) * \real{0.3279}}
  >{\raggedright\arraybackslash}p{(\linewidth - 6\tabcolsep) * \real{0.2623}}@{}}
\toprule\noalign{}
\begin{minipage}[b]{\linewidth}\raggedright
Data set (access, licence)
\end{minipage} & \begin{minipage}[b]{\linewidth}\raggedright
Outcome(s)
\end{minipage} & \begin{minipage}[b]{\linewidth}\raggedright
Predictors
\end{minipage} & \begin{minipage}[b]{\linewidth}\raggedright
Development / external split
\end{minipage} \\
\midrule\noalign{}
\endhead
\bottomrule\noalign{}
\endlastfoot
\textbf{International Stroke Trial}, IST (Edinburgh
DataShare, doi:10.7488/ds/104; ODC-By, plus the DataShare
end-user licence accepted by the author). The file has been
downloaded and has not been tabulated. &
\texttt{ist\_death14}: death within 14 days (patients with
known 14-day status); \texttt{ist\_dd6}: dead or dependent
at 6 months (unknown status excluded) & the 21 baseline
variables recorded at randomisation: age; sex; delay from
onset (h); systolic blood pressure; conscious state (alert,
drowsy, unconscious); symptoms on waking; atrial
fibrillation; aspirin within 3 days; heparin within 24 h; CT
status (not done, done without visible infarct, visible
infarct); stroke subtype (TACS, PACS, LACS, POCS, other);
the eight deficit items; and a flag for ``a deficit could
not be assessed'' & primary: UK (development) against all
other countries (external). Sensitivity: randomised before
1994 against 1994-1996. If the UK pool has fewer than 3,500
complete cases: UK, Ireland and the Northern European
countries (Norway, Sweden, Denmark, Finland, the
Netherlands). \\
\textbf{CRASH-2} (LSHTM Data Compass / freeBIRD,
doi:10.17037/DATA.10; CC BY as stated by freeBIRD). The
download waits for the author's freeBIRD registration. &
\texttt{crash2\_death28}: death in hospital within 28 days
of randomisation & following Perel et al.~(2012): Glasgow
coma score, age, systolic blood pressure, heart rate,
respiratory rate, hours since injury (continuous), injury
type (blunt only or any penetrating), and sex. The data
dictionary's missing-value codes are set to missing. &
temporal: randomised before 1 January 2008 against
2008-2010. The public file has no country. If the early pool
has fewer than 3,500 complete cases, the cut moves to the
date that splits the complete-case table in half. \\
\textbf{Sepsis survival, minimal clinical records} (UCI
repository id 827, doi:10.24432/C53C8N; CC BY 4.0) &
\texttt{sepsis\_death}: death during the admission & age,
sex, and septic episode number & primary: the 91,153
admissions not meeting Sepsis-3 (development) against the
19,051 Sepsis-3 admissions (external). Sensitivity: a random
half split. \\
\textbf{PhysioNet/CinC Challenge 2012} (PhysioNet, ODC-By
1.0) & \texttt{physionet\_ihd}: in-hospital death & 27
summaries of the first 48 hours, following the SAPS-I
variable list: age; sex; ICU type; the minimum and maximum
of heart rate, systolic blood pressure, temperature,
respiratory rate, white cell count, glucose, potassium and
sodium; the minimum of GCS, haematocrit and bicarbonate; the
maximum of BUN, creatinine and bilirubin; total urine
output; any mechanical ventilation & primary: sets A + B
(development) against set C (held out). Sensitivity:
coronary-care, cardiac-surgery and surgical ICUs against the
medical ICU. \\
\end{longtable}
}

\begin{itemize}
\tightlist
\item
  \textbf{The recalibration rows.} In every external region,
  500 random rows (fixed seed) form the local recalibration
  sample, and the remaining rows form the evaluation set.
\item
  \textbf{If CRASH-2 cannot be obtained.} The real-data
  analysis proceeds with the other three data sets, and this
  is reported.
\end{itemize}

\textbf{Missing data.}

\begin{itemize}
\tightlist
\item
  IST and CRASH-2: complete-case analysis on the predictors.
  Excluded rows are counted per data set and split. In IST,
  atrial fibrillation is missing for the 984 pilot-phase
  patients (5.1 \%).
\item
  Sepsis: no missing values.
\item
  PhysioNet 2012: complete cases would lose most stays.
  Inside each development subsample, medians are learned on
  the development rows, and a missing-value indicator is
  added for every variable missing in at least 5 \% of those
  rows. Every arm gets the same imputed matrix.
\item
  A sensitivity analysis keeps the incomplete rows (IST,
  CRASH-2) or uses the TFMs' native missing-value handling
  (PhysioNet).
\end{itemize}

\subsubsection{10.2 Learning curves}\label{learning-curves}

\begin{itemize}
\tightlist
\item
  \textbf{Subsamples.} For each data set and outcome: B =
  100 development subsamples at each n in \{200, 500, 1000,
  3000\}, drawn by simple random sampling from the
  development region. The seeds are derived from (data set,
  split, n, b).
\item
  \textbf{Evaluation.} Each subsample is evaluated on the
  \textbf{whole} external evaluation set.
\item
  \textbf{Arms.} All arms of Section 5.2 except the oracle,
  with the categorical declarations of Section 5.1.
\item
  \textbf{Small regions.} If the development region holds
  fewer than 3,000 rows, the largest n is the region size.
\item
  \textbf{Few events.} Subsamples with very few events are
  kept.
\item
  \textbf{Estimands:}

  \begin{itemize}
  \tightlist
  \item
    empirical calibration slope and
    calibration-in-the-large;
  \item
    empirical ICI, E50 and E90;
  \item
    Brier score and AUC;
  \item
    the empirical analogue of the calibration class
    (reported as empirical, not true);
  \item
    the battery of Section 5.3 on the external evaluation
    set.
  \end{itemize}
\end{itemize}

\subsubsection{10.3 Localization and repair on real
data}\label{localization-and-repair-on-real-data}

\begin{itemize}
\tightlist
\item
  \textbf{Localization.} \texttt{localize.external()} on the
  external evaluation set for FLIC and the TFMs, with these
  covariates: age band, sex, region (where present) and a
  severity score (conscious state in IST, GCS in CRASH-2,
  SAPS-I in PhysioNet).
\item
  \textbf{Repair.} The recalibration ladder of Section 5.5
  is fitted on the 500-row local sample and evaluated on the
  rest of the external region.
\end{itemize}

\subsubsection{10.4 Plasmode tier}\label{plasmode-tier}

A known true risk is placed on real covariate distributions,
with outcomes that no TFM could have seen during
pretraining.

\begin{itemize}
\tightlist
\item
  \textbf{Truth model.}

  \begin{itemize}
  \tightlist
  \item
    A main-effects logistic model is fitted once to the real
    outcome on the full complete-case table, with
    standardized design columns.
  \item
    Two terms of fixed size are then added on the logit
    scale: a U-shaped term 0.4 (z\^{}2 - 1) in one
    standardized continuous predictor, and an interaction
    0.6 z\_a x\_b.
  \item
    The intercept is solved by bisection so that the mean p*
    equals 10 \% and 30 \%.
  \end{itemize}

  {\def\LTcaptype{none} % do not increment counter
  \begin{longtable}[]{@{}
    >{\raggedright\arraybackslash}p{(\linewidth - 4\tabcolsep) * \real{0.3333}}
    >{\raggedright\arraybackslash}p{(\linewidth - 4\tabcolsep) * \real{0.3333}}
    >{\raggedright\arraybackslash}p{(\linewidth - 4\tabcolsep) * \real{0.3333}}@{}}
  \toprule\noalign{}
  \begin{minipage}[b]{\linewidth}\raggedright
  Data set
  \end{minipage} &
  \begin{minipage}[b]{\linewidth}\raggedright
  U-shaped term
  \end{minipage} &
  \begin{minipage}[b]{\linewidth}\raggedright
  Interaction
  \end{minipage} \\
  \midrule\noalign{}
  \endhead
  \bottomrule\noalign{}
  \endlastfoot
  IST (\texttt{ist\_dd6}, primary) & systolic blood pressure
  & age x impaired consciousness \\
  CRASH-2 & systolic blood pressure & GCS x penetrating
  injury \\
  PhysioNet 2012 & maximum heart rate & age x mechanical
  ventilation \\
  \end{longtable}
  }
\item
  \textbf{Design.} R = 200 replicates per n in \{200, 500,
  1000, 3000\}, with development subsamples from the
  development region and evaluation on the external
  evaluation set.
\item
  \textbf{Analyses.} The oracle estimands, the summaries of
  research questions Q1-Q3 (as secondary checks) and the
  size of the battery are computed as in the simulation.
\item
  \textbf{Frozen truths.} Each truth is stored with a
  SHA-256 of p* and checked before each use.
\end{itemize}

\subsubsection{10.5 Reporting}\label{reporting}

The real-data analyses are reported following TRIPOD+AI
(Collins et al.~2024). Real-data results are descriptive and
supportive; they do not enter the labels of Section 7.

\begin{center}\rule{0.5\linewidth}{0.5pt}\end{center}

\subsection{11. Primary, secondary and exploratory
analyses}\label{primary-secondary-and-exploratory-analyses}

\textbf{Primary.}

\begin{enumerate}
\def\labelenumi{\arabic{enumi}.}
\tightlist
\item
  The oracle estimands: the true calibration slope b*, the
  true intercept a* and the true ICI of every arm in the 48
  core scenarios, with MCSE.
\item
  The paired differences in true ICI, TFM minus FLIC, per
  core scenario: the mean paired difference with its MCSE
  and 95 \% interval, labelled better, worse, equivalent or
  inconclusive by the rule and margin of Section 7.0.
\item
  The answers to research questions Q1-Q5, by the reporting
  rules of Sections 7.1-7.5.
\item
  \textbf{The minimum-n-for-calibration table} (Section
  4.5), a planned primary output for practitioners: for each
  arm, at each event rate (3, 10 and 30 \%) and in each
  shape, the smallest development n at which the arm meets
  each calibration class in at least 80 \% of replicates,
  beside Riley's minimum n and the EPV. The pooled row of
  each event rate gives the largest of the four shapes. The
  same table with the lower 95 \% limit of the proportion,
  and at half and double the thresholds, goes to the
  supplement. The table is descriptive; its only labels are
  those of Q2.
\item
  \textbf{The robustness map} of the random-design tier: the
  regions of Section 6.4 where each TFM is better or worse
  than FLIC by more than 0.01, with the robustness check of
  Section 7.6 (Q6).
\end{enumerate}

\textbf{Secondary.}

\begin{itemize}
\tightlist
\item
  The rest of the meta-model of Section 6.4: partial
  effects, the D\_logb maps, and the detection and size
  maps.
\item
  All other oracle and empirical estimands,
  \texttt{true\_ll\_ici}, and empirical against true ICI.
\item
  Size and detection of every test:

  \begin{itemize}
  \tightlist
  \item
    the oracle nulls (the main size study);
  \item
    each arm's Y*;
  \item
    the detection curves;
  \item
    the validation-size subsamples;
  \item
    Stukel's size beside its power;
  \item
    the auxiliary size study.
  \end{itemize}
\item
  Localization: the proportion of replicates in which COV is
  named, on S3 and T3, for each arm. As a reference for each
  arm: the oracle share of the variance of logit(p*) not
  explained by a smooth function of logit(p-hat). A large
  share means that a covariate-level error truly exists for
  that arm.
\item
  The full recalibration ladder (R0-R3, m = 100 and 500) and
  the NPV of the test-guided rule.
\item
  The temperature arms: the paired effect of the softmax
  temperature on b*, a*, mean risk and true ICI, within
  TabPFN-3.5 and within TabICLv2.
\item
  The counts of failed, degenerate and not-run fits, and of
  the flags of Section 8.2.
\item
  Real data (Section 10).
\end{itemize}

\textbf{Exploratory} (labelled as such in the paper):

\begin{itemize}
\tightlist
\item
  the stress set beyond the severity ladder;
\item
  the legacy Real-TabPFN-2 arm;
\item
  prediction instability;
\item
  any analysis not listed in this document.
\end{itemize}

\textbf{Main-text figures and tables} are planned as:

\begin{itemize}
\tightlist
\item
  a design schematic;
\item
  b* against n;
\item
  oracle calibration curves;
\item
  minimum n by arm, and the minimum-n-for-calibration table;
\item
  the random-design maps: where each TFM is better or worse
  than FLIC;
\item
  detection curves with a size inset;
\item
  a recalibration slopegraph;
\item
  real-data learning curves and a localization map;
\item
  a summary-of-findings table (one row per research
  question);
\item
  a size table (every test on every arm).
\end{itemize}

\begin{center}\rule{0.5\linewidth}{0.5pt}\end{center}

\subsection{12. Deviations policy}\label{deviations-policy}

\begin{itemize}
\tightlist
\item
  \texttt{DEVIATIONS.md} is deposited empty with this
  registration, and again with each later version of the
  project archive.
\item
  \textbf{Recording a change.} Every change after the freeze
  is recorded there with: the date; what changed; why;
  whether any main-run result had been seen when the
  decision was made; and its expected effect.
\item
  \textbf{Minor deviations} are engineering changes that do
  not alter an estimand, a sample, a model setting or a
  decision rule. An example is a bug fix that is verified
  not to change stored predictions. They are listed and
  summarized in the supplement.
\item
  \textbf{Major deviations} are any change to the DGPs,
  arms, settings, replicate counts, estimands, thresholds,
  margins, research questions or reporting rules. They are
  listed in the paper. Results that depend on a major
  deviation are labelled exploratory, and the pre-registered
  analysis is also reported where it can be run.
\item
  \textbf{Bugs.} A bug found after results are seen is
  fixed, the affected parts are rerun in full, and both the
  bug and the rerun are recorded.
\item
  \textbf{Device changes.} A change of device for a scenario
  group under the rule of Section 9.3 is not a deviation,
  but every scenario's device is reported.
\item
  \textbf{``Pre-registered.''} The paper uses the word only
  for the analyses in this document, and only because the
  registration date precedes the main run.
\end{itemize}

\begin{center}\rule{0.5\linewidth}{0.5pt}\end{center}

\subsection{13. Software and
hardware}\label{software-and-hardware}

{\def\LTcaptype{none} % do not increment counter
\begin{longtable}[]{@{}
  >{\raggedright\arraybackslash}p{(\linewidth - 2\tabcolsep) * \real{0.4333}}
  >{\raggedright\arraybackslash}p{(\linewidth - 2\tabcolsep) * \real{0.5667}}@{}}
\toprule\noalign{}
\begin{minipage}[b]{\linewidth}\raggedright
Component
\end{minipage} & \begin{minipage}[b]{\linewidth}\raggedright
Version
\end{minipage} \\
\midrule\noalign{}
\endhead
\bottomrule\noalign{}
\endlastfoot
Python (both environments) & 3.11.15 \\
PyTorch (CUDA build) and CUDA runtime & 2.13.0+cu126, CUDA
12.6 (NVIDIA driver 591.86) \\
numpy, scipy, pandas, pyarrow & 2.4.6, 1.17.1, 2.3.3,
24.0.0 \\
scikit-learn, LightGBM (Python) & 1.9.1, 4.7.0 \\
tabpfn & 9.1.0 \\
autogluon.tabular (for \texttt{MitraClassifier}),
transformers, huggingface-hub & 1.6.3, 5.19.0, 1.33.0 \\
tabicl & 2.2.0 \\
R & 4.4.1 \\
ebrahim.gof & 2.10.0 (release branch, commit 19480aa; the
build is checked by the SHA-256 of its code) \\
givitiR, callr, logistf, glmnet, mgcv, pmsampsize & 1.3,
3.7.6, 1.26.1, 4.1-10, 1.9-4, 1.1.3 \\
arrow, data.table, digest, jsonlite & 23.0.1.2, 1.18.2.1,
0.6.39, 2.0.0 \\
Hardware & Windows 11 Pro; AMD Ryzen 9 3900X (12 cores, 24
threads); 32 GB RAM; NVIDIA GeForce GTX 1660 Ti 6 GB
(Turing; bfloat16 emulated) \\
\end{longtable}
}

\textbf{Model checkpoints.} Revision = the Hugging Face
commit of the repository in the project cache.

\begin{itemize}
\tightlist
\item
  TabPFN-3.5, standard (arms \texttt{tabpfn},
  \texttt{tabpfn\_t1}):
  \texttt{tabpfn-v3.5-20260909.safetensors}, repository
  \texttt{Prior-Labs/tabpfn\_3\_5}, revision
  \texttt{06bf2ba35c80a92a3b9abb436b99cf49e7a0365e}.

  \begin{itemize}
  \tightlist
  \item
    SHA-256
    \texttt{ece4d67eadfea42eb0e610df5189bea60cb7f31073d81e9c7a019b76eacf0be3}
  \end{itemize}
\item
  TabPFN-3.5-Fast (recorded, not used):
  \texttt{tabpfn-v3.5-fast-20260909.safetensors}, same
  repository and revision.

  \begin{itemize}
  \tightlist
  \item
    SHA-256
    \texttt{14e0b81285bca7a6d1ef121169f27f651f76dff49cad6c0b8779d5c836018550}
  \end{itemize}
\item
  Real-TabPFN-2 (arm \texttt{tabpfn\_legacy}):
  \texttt{tabpfn-v2-classifier-finetuned-zk73skhh.ckpt},
  repository \texttt{Prior-Labs/TabPFN-v2-clf}.

  \begin{itemize}
  \tightlist
  \item
    SHA-256
    \texttt{cf8c519c01eaf1613ee91239006d57b1c806ff5f23ac1aeb1315ba1015210e49}
  \end{itemize}
\item
  TabICLv2 (arms \texttt{tabicl}, \texttt{tabicl\_t1}):
  \texttt{tabicl-classifier-v2-20260212.ckpt}, repository
  \texttt{jingang/TabICL}, revision
  \texttt{4dcd344ece2c00be9e831fdd35bed57b5ad83e19}.

  \begin{itemize}
  \tightlist
  \item
    SHA-256
    \texttt{bdc7dbd5e4ff21f8f0456fcf90c6b7cdf72dbea960f2d05b19bec19f9b3d4ed0}
  \end{itemize}
\item
  Mitra-v2 (arm \texttt{mitra}): \texttt{model.safetensors},
  repository \texttt{autogluon/mitra-classifier-2}, revision
  \texttt{edada0d20759c58ada8c8605c25f22f6e98ea5f0}.

  \begin{itemize}
  \tightlist
  \item
    SHA-256
    \texttt{5ffab0e2cf52f61c5b7c7eb1e8542996736d1023a0212190cc09abc2119a1e09}
  \end{itemize}
\end{itemize}

At the start of the main run, all of these are written to
\texttt{results/environment.json} and checked against this
table. A mismatch stops the run.

\begin{center}\rule{0.5\linewidth}{0.5pt}\end{center}

\subsection{14. Archive and availability
plan}\label{archive-and-availability-plan}

\begin{itemize}
\tightlist
\item
  \textbf{This pre-registration} is registered on OSF
  Registries before the main run.
\item
  \textbf{Code} goes in a public GitHub repository under the
  MIT licence; \textbf{results, figures and tables} under CC
  BY 4.0. The repository holds \texttt{CITATION.cff} and
  \texttt{.zenodo.json} with the author and his ORCID.
\item
  \textbf{Zenodo}: a version DOI from a GitHub release. The
  paper cites the concept DOI.
\item
  \textbf{The ebrahim.gof 2.10.0 source} used here (commit
  19480aa) is archived with the code.
\end{itemize}

\textbf{Outputs for reuse (aim A8).}

\begin{itemize}
\tightlist
\item
  \textbf{(a) An open benchmark of calibrated and
  miscalibrated predictions.}

  \begin{itemize}
  \tightlist
  \item
    The full per-replicate prediction and test dataset of
    the simulation (the 58 scenarios and the random
    designs): for every evaluation record, p*, Y, the
    covariates where stored, and the p-hat and the
    exact-null Y* of every arm, as float32 parquet; and
    every row of the test layer.
  \item
    It is deposited on Zenodo under \textbf{CC BY 4.0} as
    its own record, with its own DOI, so that others can
    test a new calibration method or test on exactly the
    same predictions, with the truth known.
  \item
    It comes with a \textbf{loader}
    (\texttt{code/benchmark\_loader.py}: wide and long views
    of a replicate, iteration over a family, the study's own
    test rows) and a \textbf{data dictionary} of every
    column and metadata field, with the design tables
    \texttt{scenarios.csv} and \texttt{random\_designs.csv}.
  \end{itemize}
\item
  \textbf{(b) A calibration audit in one call.}

  \begin{itemize}
  \tightlist
  \item
    R: \texttt{ebrahim.gof::run.all.external()} followed by
    \texttt{localize.external()}, packaged as the recipe
    \texttt{R/audit\_calibration.R}
    (\texttt{audit\_calibration(y,\ p,\ X)}).
  \item
    Python: \texttt{code/audit\_calibration.py}, a
    command-line tool and function that takes y and p (and
    optionally X) from any model and runs the R recipe
    through a temporary file; rpy2 is not needed.
  \item
    Both are in the code archive, with an example on one
    benchmark replicate.
  \end{itemize}
\item
  \textbf{(c) A minimum-n-for-calibration table} by event
  rate for each model (Section 11, primary output 4),
  deposited as a csv beside the paper's table.
\item
  \textbf{(d) A TRIPOD+AI alignment table}: each TRIPOD+AI
  item (Collins et al.~2024) that applies to the real-data
  analyses (Section 10.5), with the section of the paper or
  supplement that answers it, or the reason it does not
  apply.
\end{itemize}

\textbf{What is never archived.}

\begin{itemize}
\tightlist
\item
  \textbf{TabPFN weights.} TabPFN-3.5 is used under the
  TABPFN-3.5 License v1.0, which allows non-commercial
  research use only and does not allow redistribution. No
  archive of this study holds any TabPFN weights. The
  archive holds instead
  \texttt{code/get\_tabpfn\_weights.py}, which prints the
  licence notice, downloads the checkpoints from Prior Labs
  after the users have accepted the licence themselves, and
  checks their SHA-256 against Section 13. The other
  checkpoints are referenced by repository, file name,
  revision and SHA-256.
\item
  \textbf{Real patient data.} They are not redistributed.
  The repository gives loaders and access instructions.
  Derived per-patient predictions on real data are deposited
  only for data sets whose licence allows redistribution
  with attribution (Sepsis: CC BY 4.0; PhysioNet 2012:
  ODC-By 1.0). For IST and CRASH-2, only aggregate results
  are deposited.
\item
  \textbf{A preprint} is posted when the results and the
  manuscript are final.
\end{itemize}

\begin{center}\rule{0.5\linewidth}{0.5pt}\end{center}

\subsection{15. References}\label{references}

\begin{itemize}
\tightlist
\item
  Austin PC, Steyerberg EW (2019). The Integrated
  Calibration Index (ICI) and related metrics for
  quantifying the calibration of logistic regression models.
  \emph{Statistics in Medicine} 38:4051-4065.
  doi:10.1002/sim.8281
\item
  Collins GS, et al.~(2024). TRIPOD+AI statement: updated
  guidance for reporting clinical prediction models that use
  regression or machine learning methods. \emph{BMJ}
  385:e078378. doi:10.1136/bmj-2023-078378
\item
  Ebrahim EK, Hussein OA, El-Kotory A (2026). {[}EDGE{]} A
  grouped calibration test for logistic regression that
  tolerates a few corrupted records. arXiv:2608.20511.
  doi:10.48550/arXiv.2608.20511
\item
  Hollmann N, Müller S, Purucker L, Krishnakumar A, Körfer
  M, Hoo SB, Schirrmeister RT, Hutter F (2025). Accurate
  predictions on small data with a tabular foundation model.
  \emph{Nature} 637:319-326. doi:10.1038/s41586-024-08328-6
\item
  Morris TP, White IR, Crowther MJ (2019). Using simulation
  studies to evaluate statistical methods. \emph{Statistics
  in Medicine} 38:2074-2102. doi:10.1002/sim.8086
\item
  Nattino G, Finazzi S, Bertolini G (2014). A new
  calibration test and a reappraisal of the calibration belt
  for the assessment of prediction models based on
  dichotomous outcomes. \emph{Statistics in Medicine}
  33:2390-2407. doi:10.1002/sim.6100
\item
  Osius G, Rojek D (1992). Normal goodness-of-fit tests for
  multinomial models with large degrees of freedom.
  \emph{Journal of the American Statistical Association}
  87:1145-1152. doi:10.1080/01621459.1992.10476271
\item
  Perel P, et al.~(2012). Predicting early death in patients
  with traumatic bleeding: development and validation of
  prognostic model. \emph{BMJ} 345:e5166.
  doi:10.1136/bmj.e5166
\item
  Prior Labs (2026). TabPFN-3.5 model card and TABPFN-3.5
  License v1.0. Hugging Face repository
  Prior-Labs/tabpfn\_3\_5 (checkpoint tabpfn-v3.5-20260909).
\item
  Puhr R, Heinze G, Nold M, Lusa L, Geroldinger A (2017).
  Firth's logistic regression with rare events: accurate
  effect estimates and predictions? \emph{Statistics in
  Medicine} 36:2302-2317. doi:10.1002/sim.7273
\item
  Qu J, Holzmüller D, Varoquaux G, Le Morvan M (2026).
  TabICLv2: A better, faster, scalable, and open tabular
  foundation model. arXiv:2602.11139.
\item
  Riley RD, et al.~(2020). Calculating the sample size
  required for developing a clinical prediction model.
  \emph{BMJ} 368:m441. doi:10.1136/bmj.m441
\item
  Riley RD, et al.~(2021). Minimum sample size for external
  validation of a clinical prediction model with a binary
  outcome. \emph{Statistics in Medicine} 40:4230-4251.
  doi:10.1002/sim.9025
\item
  Riley RD, Collins GS (2023). Stability of clinical
  prediction models developed using statistical or machine
  learning methods. \emph{Biometrical Journal} 65:2200302.
  doi:10.1002/bimj.202200302
\item
  Siepe BS, Bartoš F, Morris TP, Boulesteix A-L, Heck DW,
  Pawel S (2024). Simulation studies for methodological
  research in psychology: a standardized template for
  planning, preregistration, and reporting.
  \emph{Psychological Methods}. doi:10.1037/met0000695
\item
  Stukel TA (1988). Generalized logistic models.
  \emph{Journal of the American Statistical Association}
  83:426-431. doi:10.1080/01621459.1988.10478613
\item
  Tao Y, Zhang X, Liu X, Han B, Maddix D, et al.~(2026).
  Mitra-v2 technical report. arXiv:2609.04540.
\item
  Van Calster B, Nieboer D, Vergouwe Y, De Cock B, Pencina
  MJ, Steyerberg EW (2016). A calibration hierarchy for risk
  models was defined: from utopia to empirical data.
  \emph{Journal of Clinical Epidemiology} 74:167-176.
  doi:10.1016/j.jclinepi.2015.12.005
\end{itemize}

\textbf{Checks.}

\begin{itemize}
\tightlist
\item
  The journal DOIs were checked on Crossref: on 2026-10-06,
  and on 2026-10-07 for Osius-Rojek, Stukel, Hollmann et
  al.~and Siepe et al.
\item
  The arXiv titles and authors were checked on their arXiv
  records on 2026-10-07.
\item
  The dataset DOIs were checked on DataCite on 2026-10-06.
\end{itemize}

\begin{center}\rule{0.5\linewidth}{0.5pt}\end{center}

\subsection{Appendix A. Scenario
table}\label{appendix-a.-scenario-table}

Produced by \texttt{code/dgp.py\ build}
(\texttt{code/scenarios.csv}):

\begin{itemize}
\tightlist
\item
  kappa and alpha are solved on 10\^{}6 fixed draws;
\item
  the realized prevalence equals pi to 2e-15 in every row;
\item
  EPV = n pi / p;
\item
  the Riley minimum n comes from R \texttt{pmsampsize}
  (shrinkage 0.9, Cox-Snell R\^{}2 from the oracle AUC).
\end{itemize}

The scenario name gives the family, the event rate (pi03 = 3
\%) and n.~Shapes, links and gamma are in Sections 3.2 and
3.6.

{\def\LTcaptype{none} % do not increment counter
\begin{longtable}[]{@{}
  >{\raggedright\arraybackslash}p{(\linewidth - 14\tabcolsep) * \real{0.1067}}
  >{\raggedright\arraybackslash}p{(\linewidth - 14\tabcolsep) * \real{0.1867}}
  >{\raggedright\arraybackslash}p{(\linewidth - 14\tabcolsep) * \real{0.1067}}
  >{\raggedright\arraybackslash}p{(\linewidth - 14\tabcolsep) * \real{0.1067}}
  >{\raggedright\arraybackslash}p{(\linewidth - 14\tabcolsep) * \real{0.1067}}
  >{\raggedright\arraybackslash}p{(\linewidth - 14\tabcolsep) * \real{0.1333}}
  >{\raggedright\arraybackslash}p{(\linewidth - 14\tabcolsep) * \real{0.1067}}
  >{\raggedright\arraybackslash}p{(\linewidth - 14\tabcolsep) * \real{0.1467}}@{}}
\toprule\noalign{}
\begin{minipage}[b]{\linewidth}\raggedright
id
\end{minipage} & \begin{minipage}[b]{\linewidth}\raggedright
scenario
\end{minipage} & \begin{minipage}[b]{\linewidth}\raggedright
N\_eval
\end{minipage} & \begin{minipage}[b]{\linewidth}\raggedright
kappa
\end{minipage} & \begin{minipage}[b]{\linewidth}\raggedright
alpha
\end{minipage} & \begin{minipage}[b]{\linewidth}\raggedright
oracle AUC
\end{minipage} & \begin{minipage}[b]{\linewidth}\raggedright
EPV
\end{minipage} & \begin{minipage}[b]{\linewidth}\raggedright
Riley min n
\end{minipage} \\
\midrule\noalign{}
\endhead
\bottomrule\noalign{}
\endlastfoot
1 & S1\_pi03\_n200 & 3334 & 0.7327 & -4.4586 & 0.810 & 0.75
& 1662 \\
2 & S1\_pi03\_n500 & 3334 & 0.7327 & -4.4586 & 0.810 & 1.88
& 1662 \\
3 & S1\_pi03\_n1000 & 3334 & 0.7327 & -4.4586 & 0.810 & 3.75
& 1662 \\
4 & S1\_pi03\_n3000 & 3334 & 0.7327 & -4.4586 & 0.810 &
11.25 & 1662 \\
5 & S1\_pi10\_n200 & 1000 & 0.7327 & -3.0255 & 0.800 & 2.50
& 609 \\
6 & S1\_pi10\_n500 & 1000 & 0.7327 & -3.0255 & 0.800 & 6.25
& 609 \\
7 & S1\_pi10\_n1000 & 1000 & 0.7327 & -3.0255 & 0.800 &
12.50 & 609 \\
8 & S1\_pi10\_n3000 & 1000 & 0.7327 & -3.0255 & 0.800 &
37.50 & 609 \\
9 & S1\_pi30\_n200 & 1000 & 0.7327 & -1.3603 & 0.788 & 7.50
& 323 \\
10 & S1\_pi30\_n500 & 1000 & 0.7327 & -1.3603 & 0.788 &
18.75 & 323 \\
11 & S1\_pi30\_n1000 & 1000 & 0.7327 & -1.3603 & 0.788 &
37.50 & 323 \\
12 & S1\_pi30\_n3000 & 1000 & 0.7327 & -1.3603 & 0.788 &
112.50 & 323 \\
13 & S2\_pi03\_n200 & 3334 & 0.6230 & -4.4357 & 0.827 & 0.75
& 1454 \\
14 & S2\_pi03\_n500 & 3334 & 0.6230 & -4.4357 & 0.827 & 1.88
& 1454 \\
15 & S2\_pi03\_n1000 & 3334 & 0.6230 & -4.4357 & 0.827 &
3.75 & 1454 \\
16 & S2\_pi03\_n3000 & 3334 & 0.6230 & -4.4357 & 0.827 &
11.25 & 1454 \\
17 & S2\_pi10\_n200 & 1000 & 0.6230 & -2.9203 & 0.800 & 2.50
& 609 \\
18 & S2\_pi10\_n500 & 1000 & 0.6230 & -2.9203 & 0.800 & 6.25
& 609 \\
19 & S2\_pi10\_n1000 & 1000 & 0.6230 & -2.9203 & 0.800 &
12.50 & 609 \\
20 & S2\_pi10\_n3000 & 1000 & 0.6230 & -2.9203 & 0.800 &
37.50 & 609 \\
21 & S2\_pi30\_n200 & 1000 & 0.6230 & -1.2353 & 0.769 & 7.50
& 349 \\
22 & S2\_pi30\_n500 & 1000 & 0.6230 & -1.2353 & 0.769 &
18.75 & 349 \\
23 & S2\_pi30\_n1000 & 1000 & 0.6230 & -1.2353 & 0.769 &
37.50 & 349 \\
24 & S2\_pi30\_n3000 & 1000 & 0.6230 & -1.2353 & 0.769 &
112.50 & 349 \\
25 & S3\_pi03\_n200 & 3334 & 0.6142 & -4.4168 & 0.814 & 0.75
& 1612 \\
26 & S3\_pi03\_n500 & 3334 & 0.6142 & -4.4168 & 0.814 & 1.88
& 1612 \\
27 & S3\_pi03\_n1000 & 3334 & 0.6142 & -4.4168 & 0.814 &
3.75 & 1612 \\
28 & S3\_pi03\_n3000 & 3334 & 0.6142 & -4.4168 & 0.814 &
11.25 & 1612 \\
29 & S3\_pi10\_n200 & 1000 & 0.6142 & -2.9651 & 0.800 & 2.50
& 609 \\
30 & S3\_pi10\_n500 & 1000 & 0.6142 & -2.9651 & 0.800 & 6.25
& 609 \\
31 & S3\_pi10\_n1000 & 1000 & 0.6142 & -2.9651 & 0.800 &
12.50 & 609 \\
32 & S3\_pi10\_n3000 & 1000 & 0.6142 & -2.9651 & 0.800 &
37.50 & 609 \\
33 & S3\_pi30\_n200 & 1000 & 0.6142 & -1.2956 & 0.783 & 7.50
& 323 \\
34 & S3\_pi30\_n500 & 1000 & 0.6142 & -1.2956 & 0.783 &
18.75 & 323 \\
35 & S3\_pi30\_n1000 & 1000 & 0.6142 & -1.2956 & 0.783 &
37.50 & 323 \\
36 & S3\_pi30\_n3000 & 1000 & 0.6142 & -1.2956 & 0.783 &
112.50 & 323 \\
37 & S4\_pi03\_n200 & 3334 & 0.6697 & -4.3699 & 0.798 & 0.75
& 1845 \\
38 & S4\_pi03\_n500 & 3334 & 0.6697 & -4.3699 & 0.798 & 1.88
& 1845 \\
39 & S4\_pi03\_n1000 & 3334 & 0.6697 & -4.3699 & 0.798 &
3.75 & 1845 \\
40 & S4\_pi03\_n3000 & 3334 & 0.6697 & -4.3699 & 0.798 &
11.25 & 1845 \\
41 & S4\_pi10\_n200 & 1000 & 0.6697 & -3.0419 & 0.800 & 2.50
& 609 \\
42 & S4\_pi10\_n500 & 1000 & 0.6697 & -3.0419 & 0.800 & 6.25
& 609 \\
43 & S4\_pi10\_n1000 & 1000 & 0.6697 & -3.0419 & 0.800 &
12.50 & 609 \\
44 & S4\_pi10\_n3000 & 1000 & 0.6697 & -3.0419 & 0.800 &
37.50 & 609 \\
45 & S4\_pi30\_n200 & 1000 & 0.6697 & -1.5964 & 0.814 & 7.50
& 323 \\
46 & S4\_pi30\_n500 & 1000 & 0.6697 & -1.5964 & 0.814 &
18.75 & 323 \\
47 & S4\_pi30\_n1000 & 1000 & 0.6697 & -1.5964 & 0.814 &
37.50 & 323 \\
48 & S4\_pi30\_n3000 & 1000 & 0.6697 & -1.5964 & 0.814 &
112.50 & 323 \\
49 & T1\_pi10\_n500 & 1000 & 0.8205 & -3.0442 & 0.800 & 6.25
& 609 \\
50 & T1\_pi10\_n3000 & 1000 & 0.8205 & -3.0442 & 0.800 &
37.50 & 609 \\
51 & T2\_pi10\_n500 & 1000 & 0.7327 & -3.0255 & 0.800 & 1.67
& 2284 \\
52 & T2\_pi10\_n3000 & 1000 & 0.7327 & -3.0255 & 0.800 &
10.00 & 2284 \\
53 & T3\_pi10\_n500 & 1000 & 0.5992 & -2.9981 & 0.800 & 4.17
& 914 \\
54 & T3\_pi10\_n3000 & 1000 & 0.5992 & -2.9981 & 0.800 &
25.00 & 914 \\
55 & T4a\_pi10\_n500 & 1000 & 0.6762 & -2.9646 & 0.800 &
6.25 & 609 \\
56 & T4a\_pi10\_n3000 & 1000 & 0.6762 & -2.9646 & 0.800 &
37.50 & 609 \\
57 & T4b\_pi10\_n500 & 1000 & 0.5297 & -2.8686 & 0.800 &
6.25 & 609 \\
58 & T4b\_pi10\_n3000 & 1000 & 0.5297 & -2.8686 & 0.800 &
37.50 & 609 \\
\end{longtable}
}

\subsection{Appendix B. The random
designs}\label{appendix-b.-the-random-designs}

Produced by \texttt{code/dgp\_random.py\ build\ -\/-D\ 1000}
(\texttt{code/random\_designs.csv}, one row per design with
its full specification). Every design was solved: the
realized prevalence equals pi to 6e-14 and the oracle AUC
equals its target to 2e-10, on 200,000 calibration draws per
design.

{\def\LTcaptype{none} % do not increment counter
\begin{longtable}[]{@{}
  >{\raggedright\arraybackslash}p{(\linewidth - 10\tabcolsep) * \real{0.3333}}
  >{\raggedright\arraybackslash}p{(\linewidth - 10\tabcolsep) * \real{0.1333}}
  >{\raggedright\arraybackslash}p{(\linewidth - 10\tabcolsep) * \real{0.1333}}
  >{\raggedright\arraybackslash}p{(\linewidth - 10\tabcolsep) * \real{0.1333}}
  >{\raggedright\arraybackslash}p{(\linewidth - 10\tabcolsep) * \real{0.1333}}
  >{\raggedright\arraybackslash}p{(\linewidth - 10\tabcolsep) * \real{0.1333}}@{}}
\toprule\noalign{}
\begin{minipage}[b]{\linewidth}\raggedright
Factor
\end{minipage} & \begin{minipage}[b]{\linewidth}\raggedright
min
\end{minipage} & \begin{minipage}[b]{\linewidth}\raggedright
10 \%
\end{minipage} & \begin{minipage}[b]{\linewidth}\raggedright
median
\end{minipage} & \begin{minipage}[b]{\linewidth}\raggedright
90 \%
\end{minipage} & \begin{minipage}[b]{\linewidth}\raggedright
max
\end{minipage} \\
\midrule\noalign{}
\endhead
\bottomrule\noalign{}
\endlastfoot
event rate pi & 0.020 & 0.027 & 0.088 & 0.290 & 0.400 \\
development n & 150 & 200 & 738 & 2732 & 3993 \\
predictors p & 4 & 6 & 17 & 28 & 30 \\
correlation rho & 0.00 & 0.06 & 0.29 & 0.53 & 0.60 \\
oracle AUC & 0.650 & 0.673 & 0.774 & 0.870 & 0.900 \\
w\_nl & 0.00 & 0.10 & 0.51 & 0.91 & 1.00 \\
w\_int & 0.00 & 0.11 & 0.51 & 0.92 & 1.00 \\
w\_link & 0.00 & 0.10 & 0.50 & 0.90 & 1.00 \\
discrete fraction & 0.00 & 0.05 & 0.24 & 0.44 & 0.50 \\
N\_eval & 1000 & 1000 & 1130 & 3730 & 4996 \\
expected development events (n pi) & 3.2 & 12.6 & 67.6 &
369.5 & 1255.4 \\
EPV & 0.12 & 0.59 & 4.01 & 26.11 & 187.37 \\
share of p* \textgreater= 0.5 & 0.000 & 0.000 & 0.014 &
0.165 & 0.337 \\
\end{longtable}
}

\begin{itemize}
\tightlist
\item
  The 1,000 designs hold 3,111 binary and 1,052 categorical
  predictors; 594 designs have at least one categorical
  predictor.
\item
  658 designs have p \textgreater{} 12, so DeepGOF-1 uses
  its axes reading there.
\item
  TabPFN-3.5 runs on the GPU in 821 designs and on a CPU
  worker in 179.
\item
  Mitra-v2 is recorded as ``does not fit in 6 GB GPU
  memory'' in 170 designs.
\item
  In 722 designs n is below Riley's minimum n, and 67
  designs expect fewer than 10 development events.
\end{itemize}

\end{document}
